\section{Introduction to Electroacoustics}

\subsection{Electrodynamic Loudspeaker: Structure}

When a voltage is applied to the terminals of a loudspeaker, an electrical current flows through its coil. This current interacts with the magnetic field of the loudspeaker, generating a mechanical force that moves the membrane (diaphragm). The membrane’s motion is sinusoidal and propagates sound waves in air.

\begin{figure}[H]
    \centering
    \includegraphics[width=0.45\textwidth]{loudspeaker_cutaway.png}
    \hfill
    \includegraphics[width=0.45\textwidth]{Loudspeaker_Structure.png}
    \caption{Simplified illustration of the electrodynamic loudspeaker. The coil is immersed in a magnetic field created by the permanent magnet. The suspension and cone convert the electrical input into acoustic motion.}
\end{figure}

In a typical configuration:
\begin{itemize}
    \item The \textbf{north pole} of the magnet is located at the top and the \textbf{south pole} at the bottom.
    \item The \textbf{voice coil} is positioned in the magnetic gap.
    \item The \textbf{suspension} (or spider) stabilizes and regulates the diaphragm’s movement.
\end{itemize}

We refer to such systems as \textbf{electromagnetic loudspeakers} because they rely on the interaction between an electric current and a magnetic field.
\[
\text{Electricity} \Rightarrow \text{Sound: Loudspeaker}
\qquad \qquad \qquad 
\text{Sound} \Rightarrow \text{Electricity: Microphone}
\]
An \textbf{electrodynamic microphone} has a structure similar to that of the loudspeaker, but operates inversely: acoustic pressure moves the diaphragm, which induces an electrical current in the coil.

\subsubsection{Magnetic Interaction and Lorentz Force}

A moving electric charge \( q \) in a magnetic field \( \underline{B} \) (expressed in Teslas, \( [T] = \frac{N \cdot s}{C \cdot m} \)) experiences a force \( \underline{F} \) given by the \textbf{Lorentz law}:
\[
\underline{F} = q \cdot \underline{u} \times \underline{B}
\]
where \( \underline{u} \) is the velocity of the particle. The force is always orthogonal to both \( \underline{u} \) and \( \underline{B} \).  
Its magnitude is expressed as:
\[
|\underline{F}| = q\,|\underline{u}|\,|\underline{B}| \sin(\alpha)
\]
where \( \alpha \) is the angle between \( \underline{u} \) and \( \underline{B} \).

\begin{figure}[H]
    \centering
    \includegraphics[width=0.5\textwidth]{lorentz_force.png}
    \caption{Representation of the Lorentz force acting on a charge moving within a magnetic field.}
\end{figure}

\subsubsection{Force on a Conductor}

In a conductor, many electrons move simultaneously. Each infinitesimal charge \( dq \) is subjected to a small force \( d\underline{F} \). Summing over all these contributions gives the total force acting on the wire.

For an infinitesimal volume element \( dV \), we can write:
\[
d\underline{F} = (-q\,\underline{u} \times \underline{B})\,n\,dV
\]
where: \( n \) is the electron concentration \( [\text{cm}^{-3}] \), \( -q \) represents the charge of a single electron.


\begin{figure}[H]
    \centering
    \includegraphics[width=0.5\textwidth]{wire_force.png}
    \caption{Infinitesimal element of wire subject to Lorentz force.}
\end{figure}



If \( dV = S \, dl \) is the infinitesimal volume of the conductor (with \( S \) its cross-section and \( dl \) an infinitesimal length), the elementary force can also be written as:
\[
d\underline{F} = (-q\,\underline{u} \times \underline{B})\,n\,S\,dl
\]
Integrating along the wire, we obtain:
\[
|\underline{F}| = \int_0^L |d\underline{F}| = nq|\underline{u}||\underline{B}|\sin(\alpha)\,S \int_0^L dl = nq|\underline{u}||\underline{B}|\sin(\alpha)\,S\,L
\]
By definition, the electric current is the rate of charge flow:
\[
I = \frac{\Delta Q}{\Delta t}
\]
Setting $\Delta t = 1\,\mathrm{s}$, we have $\Delta Q = q \, n \, dV$, thus:
\[
I = \Delta Q = q \, n \, S \, l
\]
Since $l = |\underline{u}| \, \Delta t$ and $\Delta t = 1$, we can write $l = |\underline{u}|$.

Therefore, the total force becomes:
\[
|\underline{F}| = q\,n\,|\underline{u}|\,S\,L\,|\underline{B}|\,\sin(\alpha)
\Rightarrow |\underline{F}| = I\,|\underline{B}|\,l\,\sin(\alpha)
\]

\textbf{Remark:} In loudspeakers, the coil is placed such that $\alpha = \frac{\pi}{2}$, ensuring that $\sin(\alpha)=1$ and the magnetic force is maximized.

\subsubsection{Electrical–Mechanical–Acoustical Conversion}

\begin{figure}[H]
    \centering
    \begin{minipage}[c]{0.6\linewidth}
        In a loudspeaker, the current flowing in the coil generates a magnetic field, which interacts with the permanent magnet and exerts a force on the diaphragm. This motion converts electrical energy into mechanical energy, and then into acoustic energy.
        
        \medskip
        The conversion chain is as follows:
            \[
            \text{Electrical Power } P_E \longrightarrow \text{Mechanical Power } P_M \longrightarrow \text{Acoustic Power } P_A
            \]
            The electrical power is defined as:
            \[
            P_E = V \, I
            \]
            In acoustics, we introduce the \textbf{volume velocity}:
\[
V = \underline{u} \cdot S
\]
where $S$ is the diaphragm surface area and $\underline{u}$ its particle velocity.

The acoustic power is then:
\[
P_A = P_M \, V = \frac{F}{S} \, \underline{u} \, S = F \, \underline{u}
\]

In an ideal system, all mechanical power is converted into acoustic power, i.e.
\[
P_M = P_A
\]
    \end{minipage}\hfill
    \begin{minipage}[c]{0.35\linewidth}
        \centering
        \includegraphics[width=\linewidth]{loudspeaker_force_convertion.png}
    \caption{Force generation in the loudspeaker coil: the interaction of current and magnetic field produces diaphragm motion.}
    \end{minipage}
\end{figure}



\subsection{How to Approach the Loudspeaker Description}

Before modeling a loudspeaker, it is essential to clearly understand the various components and their physical interactions. The analysis proceeds according to the following steps:

\begin{enumerate}
    \item Identify all the components involved in the system.
    \item Understand the physical limitations and possible damages.
    \item Determine where and how the Lorentz force acts.
    \item Refine parameters for good loudspeaker design (useful for tools such as Loudspeaker Design Software, LDS).
    \item Establish the relationships between sound and current.
    \item Define the fundamental sound generation equation.
    \item Combine all these aspects into a unified model.
\end{enumerate}

For \textbf{woofers} (low-frequency drivers), a larger diaphragm is preferred, as it allows efficient reproduction of low frequencies. Conversely, for \textbf{tweeters} (high-frequency drivers), smaller diaphragms perform better due to their lower mass and faster response.

\subsection{Electrodynamic Loudspeaker: Model}

The loudspeaker receives as input a pair \((V, I)\) corresponding to voltage and current applied to the \textbf{voice coil}. This electrical signal is converted into mechanical motion through electromagnetic interaction.

The energy conversion can be viewed as a succession of domains:
\[
\text{Electrical Domain: } (V, I) \quad \xrightarrow{\text{conversion}} \quad \text{Mechanical Domain: } (\underline{F}, \underline{u})
\]
\clearpage
Here:
\begin{itemize}
    \item The \textbf{electrical domain} variables are voltage \(V\) and current \(I\).
    \item The \textbf{mechanical domain} variables are force \(\underline{F}\) and velocity \(\underline{u}\).
\end{itemize}

\noindent The diaphragm suspension can be modeled mechanically using mass–spring–damper analogies, which capture the elastic and dissipative behavior of the suspension system.

From the mechanical domain, we then move to the \textbf{acoustical domain}, where motion of the diaphragm interacts with air:
\[
\text{Acoustical Domain: } (P, \underline{U})
\]
where \(P\) is the acoustic pressure and \(\underline{U}\) is the volume velocity (product of particle velocity and surface area).

\begin{figure}[H]
    \centering
    \includegraphics[width=0.75\textwidth]{conversion_chain.png}
    \caption{Block diagram showing the energy conversion from electrical to mechanical to acoustical domains.}
\end{figure}


Ultimately, air and the surrounding environment form the coupling medium that transmits acoustic power away from the loudspeaker.

\subsection{Electrical Domain Components}

The loudspeaker’s electrical behavior is described by its impedance, which links voltage and current through Ohm’s law in the frequency domain:
\[
Z_e = \frac{V_e}{i} \Rightarrow i(t) = \frac{1}{Z_e} V_e(t)
\]
\noindent\textbf{Resistor:} For a resistor of value $R$,
\[
V_R = R \, i
\]
\noindent\textbf{Capacitor:} For a capacitance $C$, the voltage–charge relationship is $V_C = \frac{1}{C} Q$, where $Q$ is the charge such that $\frac{dQ}{dt} = i$. Therefore:
\[
V_C = \frac{1}{C} \int_0^t i(\tau) \, d\tau \quad \Leftrightarrow \quad i_C(t) = C \frac{dV_C}{dt}
\]

\noindent\textbf{Inductor:} For an inductance $L$,
\[
V_L = L \frac{di}{dt} \Rightarrow i(t) = \frac{1}{L} \int_0^t V_L(\tau) \, d\tau
\]

\begin{figure}[H]
    \centering
    \includegraphics[width=0.85\textwidth]{electrical_elements.png}
    \caption{Electrical domain components: resistor, capacitor, and inductor in a loudspeaker circuit.}
\end{figure}


Each of these elements will later find its mechanical or acoustical analog in the modeling chain.


\subsection{Mechanical Domain Components}

The mechanical domain describes how the diaphragm and suspension react to forces. The three main components are mass, spring, and damper, which correspond respectively to mechanical inertia, elasticity, and losses.

\subsubsection{Mass in the Mechanical Domain}

We denote the diaphragm’s mass by $M_M$. When a force $\underline{F}$ is applied, the diaphragm acquires a velocity $\underline{u}$.

From Newton’s second law:
\[
\underline{a} = \frac{1}{M_M}\underline{F} \quad \Rightarrow \quad \underline{F} = M_M \underline{a} = M_M \frac{d\underline{u}}{dt}
\]
By integration:
\[
\underline{u}(t) = \frac{1}{M_M} \int_0^t \underline{F}(\tau)\, d\tau
\]

\noindent The element $M_M$ is modeled as a \textbf{mechanical inductance}, since the force depends on the derivative of velocity (analogous to voltage depending on the derivative of current in an inductor).

\subsubsection{Spring in the Mechanical Domain}

For a spring, the displacement $x$ is proportional to the applied force:
\[
x = C_M \, F
\]
where $C_M$ is the \textbf{mechanical compliance}, the inverse of stiffness ($K = 1/C_M$). We have:
\[
F = \frac{1}{C_M} x = K \, x 
\quad \Rightarrow \quad 
x = \int_0^t \underline{u}(\tau) d\tau = \frac{F}{K}
\]
Differentiating with respect to time:
\[
\underline{u}(t) = \frac{1}{K} \frac{d\underline{F}}{dt}
\]
Thus, a spring behaves like a \textbf{mechanical capacitance}, where displacement stores potential energy.


\subsubsection*{Damper in the Mechanical Domain}

The damper represents friction or air resistance. Its constitutive relation is linear:
\[
\underline{u} = \frac{1}{R_M} \underline{F}
\]
where $R_M$ is the \textbf{mechanical resistance}. The damper dissipates energy, analogous to an electrical resistor.

\begin{figure}[H]
    \centering
    \includegraphics[width=0.8\textwidth]{mechanical_blocks.png}
    \caption{Mechanical damper model: $\underline{u} = \frac{1}{R_M} \underline{F}$.}
\end{figure}

\clearpage 

\subsection{Analogy Theory}

When modeling electromechanical systems, it is often convenient to establish analogies between electrical and mechanical quantities. Two main types of analogies exist: the \textbf{impedance analogy} and the \textbf{admittance analogy}. These mappings allow us to represent mechanical systems using equivalent electrical circuits.


\begin{figure}[H]
    \centering
    \begin{minipage}[c]{0.45\linewidth}
        \textbf{Type 1: Impedance Analogy}

In this analogy, we map:
\[
\begin{cases}
{F} \leftrightarrow v, \\
{u} \leftrightarrow i,
\end{cases}
\]
which is particularly suitable for the impedance analogy of the mass.

\begin{itemize}
    \item Mechanical \textbf{force} corresponds to electrical \textbf{voltage}.
    \item Mechanical \textbf{velocity} corresponds to electrical \textbf{current}.
\end{itemize}
    \end{minipage}\hfill
    \begin{minipage}[c]{0.45\linewidth}
        \textbf{Type 2: Admittance Analogy}
        
        In this case, we map:
\[
\begin{cases}
{F} \leftrightarrow i, \\
{u} \leftrightarrow v,
\end{cases}
\]
which is better suited for representing the spring (capacitance analogy).
\begin{itemize}
    \item Mechanical \textbf{force} corresponds to electrical \textbf{current}.
    \item Mechanical \textbf{velocity} corresponds to electrical \textbf{voltage}.
\end{itemize}
    \end{minipage}
\end{figure}


\begin{table}[H]
\centering
\renewcommand{\arraystretch}{1.2}
\begin{tabular}{|c|c|c|}
\hline
\textbf{Mechanical Element} & \textbf{Impedance Analogy (Type 1)} & \textbf{Admittance Analogy (Type 2)} \\
\hline
Mass & Inductor & Capacitor \\
Resistor (Damper) & Resistor & Conductor (inverse of resistor) \\
Spring & Capacitor & Inductor \\
\hline
\end{tabular}
\caption{Correspondences between electrical and mechanical elements for impedance and admittance analogies.}
\end{table}

\begin{figure}[H]
    \centering
    \includegraphics[width=0.85\textwidth]{impedance_admittance_analogies.png}
    \caption{Mechanical–electrical analogies: impedance (left) and admittance (right).}
\end{figure}


\noindent\textbf{Remark:} Depending on the chosen analogy, mechanical components are represented by different electrical equivalents. Both analogies are valid and can be used based on modeling convenience.

\subsubsection{Impedance Analogy of the Spring}

Using the impedance analogy: \( {F} \leftrightarrow v, \quad 
{u} \leftrightarrow i\)

We have in the electrical domain:
\[
i = C_E \frac{dv}{dt} \Rightarrow v(t) = \frac{1}{C_E} \int_0^t i(\tau)\, d\tau
\]
By analogy, in the mechanical domain:
\[
{F}(t) = \frac{1}{C_E} \int_0^t {u}(\tau)\, d\tau
\]

\subsubsection{Admittance Analogy of the Spring}

Using the admittance analogy: $ \underline{F} \leftrightarrow i, \quad
\underline{u} \leftrightarrow v$

We obtain in the electrical domain:
\[
v = L_E \frac{di}{dt}
\]
and therefore in the mechanical domain:
\[
{u}(t) = L_E \frac{d{F}}{dt}
\]
This means that the spring can be modeled as an inductor in the admittance analogy, whereas it was represented by a capacitor in the impedance analogy.

\subsubsection{Impedance Analogy of the Damper}

Using the impedance analogy: $ \underline{F} \leftrightarrow v,\quad {u} \leftrightarrow i $

We obtain:
\[
{F} = R_M {u} \quad \Rightarrow \quad {u} = \frac{1}{R_M} {F}
\]
In the electrical domain, this becomes:
\[
v = R_E\, i \Rightarrow i(t) = \frac{1}{R_E} v(t)
\]
which corresponds to the model of an electrical resistor.

\subsubsection{Admittance Analogy of the Damper}

Using the admittance analogy: $ {F} \leftrightarrow i,\quad {u} \leftrightarrow v$

We obtain:
\[
i = G_E\, v \Rightarrow v(t) = \frac{1}{G_E} i(t)
\]
where $G_E$ is the electrical conductance (inverse of resistance). This model is equivalent to that of a resistor but expressed using admittance rather than impedance.


\subsubsection*{Impedance Analogy of the Mass}

The infinitesimal work performed by a force is:
\[
dL = F \, dx
\]
The instantaneous power in the mechanical domain is then:
\[
P_M = \frac{dL}{dt} = F \, \frac{dx}{dt} = F \, u
\]
This expression is directly analogous to the electrical power:
\[
P_E = v \, i
\]
Hence, a mechanical mass can be modeled using an \textbf{inductor} in the impedance analogy. The relationship between force and velocity mirrors that between voltage and current in an inductor:
\[
F = M_M \frac{du}{dt}
\]
The \textbf{mechano-electrical analogy} simplifies analysis by allowing us to use electrical circuit tools to model mechanical systems, avoiding complex differential equations.


\subsection{Impedances from Electrical to Mechanical Domain}



\begin{figure}[H]
    \centering
    \begin{minipage}[c]{0.55\linewidth}
        In the electrical domain, impedances are defined as the ratio between the “across” and “through” variables:
\[
Z_E = \frac{v}{i}
\]
Similarly, in the mechanical domain, we define a mechanical impedance as:
\[
Z_M = \frac{F}{u}
\]
This quantity expresses how much force is required to produce a certain velocity, analogous to voltage and current in electrical systems.
    \end{minipage}\hfill
    \begin{minipage}[c]{0.4\linewidth}
        \centering
        \includegraphics[width=\linewidth]{mechanical_impedance_analogies.png}
    \caption{Electrical domain impedance (left) and mechanical impedance (right) analogies.}
    \end{minipage}
\end{figure}




\begin{table}[H]
\centering
\renewcommand{\arraystretch}{1.3}
\begin{tabular}{|c|c|c|}
\hline
\textbf{Analogy Type} & \textbf{Mapping} & \textbf{Impedance Formula} \\
\hline
Impedance analogy & $F \leftrightarrow v$, $u \leftrightarrow i$ & $Z_M = \dfrac{F}{u} = \dfrac{\text{across}}{\text{through}}$ \\
\hline
Admittance analogy & $F \leftrightarrow i$, $u \leftrightarrow v$ & $Z_M = \dfrac{F}{u} = \dfrac{\text{through}}{\text{across}}$ \\
\hline
\end{tabular}
\caption{Mechanical impedance under impedance and admittance analogies.}
\end{table}


\noindent In both analogies, the definition of mechanical impedance remains consistent. However, the formulation of equivalent circuit relationships such as voltage dividers will differ according to the mapping used.

In the admittance analogy, the traditional “voltage divider” formula is modified.

\begin{figure}[H]
    \centering
    \includegraphics[width=0.75\textwidth]{impedance_admittance_parralel.png}\caption{Mechanical impedance combinations under impedance (left) and admittance (right) analogies.}
\end{figure}



With the impedance analogy, we have:
\[
Z_M = \frac{F}{u} = \frac{F}{u_1 + u_2} = \frac{F}{\frac{F}{Z_{M1}} + \frac{F}{Z_{M2}}}
\quad \Rightarrow \quad
Z_M = \frac{Z_{M1} Z_{M2}}{Z_{M1} + Z_{M2}}
\]
This is analogous to \textbf{impedances connected in parallel} in the electrical domain.

In the admittance analogy, we have instead:
\[
Z_M = \frac{F}{u} = \frac{F_1 + F_2}{u} = \frac{F_1}{u} + \frac{F_2}{u} 
\quad \Rightarrow \quad
Z_M = Z_{M1} + Z_{M2}
\]
which is analogous to the \textbf{parallel of admittances}.


\subsection{Acoustical Domain Components}

In the acoustical domain, we work with the quantities \( P \) and \( U \), corresponding respectively to \textbf{pressure} and \textbf{volume velocity}.  
The volume velocity is defined as:
\[
V = u \, S \quad \left[ \frac{m^3}{s} \right]
\]
where \( u \) is the particle velocity and \( S \) the surface area through which air moves.


\subsubsection{Acoustic Mass}

We model here the inertia of air moving in a cylindrical volume. The force acting on the air mass is given by Newton’s law:
\[
F = M_M \frac{du}{dt}
\]
where \( M_M \) is the mass of the volume of gas under consideration.

Since \( P = \frac{F}{S} \), we can write:
\[
P = \frac{M_M}{S} \frac{du}{dt} = \frac{M_M}{S^2} \frac{dU}{dt} = M_A \frac{dU}{dt}
\]
where:
\[
M_A = \frac{M_M}{S^2} \quad \left[ \frac{kg}{m^4} \right]
\]
is the \textbf{acoustic mass}, representing the uncompressed mass of air moving at velocity \( u \) through a surface \( S \).

\subsubsection{Impedance Analogy of the Acoustic Mass}


\begin{figure}[H]
    \centering
    \begin{minipage}[c]{0.7\linewidth}
        Using the impedance analogy: $ P \leftrightarrow v,\quad
U \leftrightarrow i$ an acoustic mass behaves as an \textbf{inductor} with inductance \( L = M_A \).

In the electrical domain, the inductor’s equation is:
\[
v = L_E \frac{di}{dt}
\]
and by analogy, in the acoustic domain:
\[
P = M_A \frac{dU}{dt} \equiv L_E \frac{di}{dt}
\]
    \end{minipage}\hfill
    \begin{minipage}[c]{0.25\linewidth}
        \centering
        \includegraphics[width=\linewidth]{acoustic_mass_impedance.png}
    \caption{Model of the mass for the impedance analogy.}
    \end{minipage}
\end{figure}

\subsubsection{Admittance Analogy of the Acoustic Mass}


\begin{figure}[H]
    \centering
    \begin{minipage}[c]{0.7\linewidth}
Using the admittance analogy: $  P \leftrightarrow i,\quad U \leftrightarrow v $
an acoustic mass can instead be represented by a \textbf{capacitor} with capacitance \( C_E = M_A \).

In the electrical domain:
\[
i = C_E \frac{dv}{dt}
\]
and in the acoustic domain:
\[
P = M_A \frac{dU}{dt} = C_E \frac{dU}{dt}
\]
    \end{minipage}\hfill
    \begin{minipage}[c]{0.25\linewidth}
        \centering
        \includegraphics[width=\linewidth]{acoustic_mass_admittance.png}
    \caption{Model of the mass for the impedance analogy.}
    \end{minipage}
\end{figure}


\clearpage 

\subsubsection{Acoustic Resistance}

We consider a cylindrical volume of gas (e.g., air) with radius \( r \).  
In the mechanical domain, the law of dampers gives:
\[
F = R_M \, u
\]
Dividing both sides by the surface \( S \), we obtain:
\[
\frac{F}{S} = \frac{R_M}{S^2} \, U \quad \Rightarrow \quad R_A = \frac{R_M}{S^2} \quad \left[ \frac{N \, s}{m^5} \right]
\]
where \( R_A \) is the \textbf{acoustic resistance}.  
Thus, a volume of air can be represented by a resistor in the electrical domain, corresponding to viscous friction losses in air motion.

We now consider a cube of height \( 2b \), width \( w \), and depth \( t \), drilled with holes of radius \( a \) that pass through the entire volume.  
Each hole can be modeled as a combination of:
\begin{itemize}
    \item an \textbf{acoustic resistance} (due to viscosity and friction), and  
    \item an \textbf{acoustic mass} (due to the inertia of air moving through the hole).
\end{itemize}

Under the impedance analogy, these two elements are in series, hence:
\[
Z_A = R_A + M_A
\]
The analytical expressions for these quantities are:
\[
R_A = \frac{1}{\pi a^2} \sqrt{\frac{2 \omega \mu}{\rho_0}} \cdot \left[ \frac{t}{a} + 2 \left( 1 - \frac{\pi a^2}{b^2} \right) \right]
, \qquad 
M_A = \frac{\rho_0}{\pi a^2} \left[ t + 17 a \left( 1 - \frac{a}{b} \right) \right]
\]
where
\begin{itemize}
    \item \( \mu \) is the \textbf{dynamic viscosity} (friction coefficient),
    \item \( \rho_0 \) is the \textbf{density of air}.
\end{itemize}

\subsubsection{Acoustic Compliance}

Air behaves like a spring. Consider a box with a movable wall, where air compression and rarefaction occur as the wall moves. The displacement is given by:
\[
x = C_M F
\]
where \( C_M \) is the \textbf{mechanical compliance}. 
\begin{figure}[H]
    \centering
    \includegraphics[width=0.35\textwidth]{acoustic_compliance_box.png}\caption{Acoustic compliance: air in a closed box behaves like a spring.}
\end{figure}
Differentiating, we obtain:
\[
u(t) = C_M \frac{dF}{dt} \quad \Leftrightarrow \quad F(t) = \frac{1}{C_M} \int_0^t u(\tau) d\tau
\]
Considering the acoustic pressure \( P = \frac{F}{S} \) as the variable of interest:
\[
P = \frac{F(t)}{S} = \frac{1}{S C_M} \int_0^t u(\tau) d\tau = \frac{1}{S^2 C_M} \int_0^t U(\tau) d\tau = \frac{1}{C_A} \int_0^t U(\tau) d\tau
\]
where
\[
C_A = C_M S^2 \quad \left[ \frac{m^5}{N} \right]
\]
is the \textbf{acoustic compliance}, representing the ability of air to store potential (compressive) energy




\subsubsection{Impedance Analogy of the Acoustic Compliance}

Using the impedance analogy: $ 
P \leftrightarrow v,\quad U \leftrightarrow i$ the equation
\[
P = \frac{1}{C_A} \int_0^t U(\tau) d\tau
\]
corresponds to the capacitor relation in the electrical domain:
\[
v(t) = \frac{1}{C_E} \int_0^t i(\tau) d\tau
\]
Thus, a volume of air can be modeled by a \textbf{capacitor} using the impedance analogy.

\subsubsection{Admittance Analogy of the Acoustic Compliance}

Using the admittance analogy: $ P \leftrightarrow i,\quad U \leftrightarrow v $ the same relation becomes:
\[
P = \frac{1}{C_A} \int_0^t U(\tau) d\tau \quad \Leftrightarrow \quad U(t) = C_A \frac{dP}{dt}
\]
which corresponds to the inductor equation:
\[
v = L_E \frac{di}{dt}
\]
Hence, a volume of air can be modeled by an \textbf{inductor} using the admittance analogy.



\subsubsection{Specific Acoustic Impedance}

We consider an acoustical element represented by a dipole with impedance:
\[
Z_A = \frac{P}{U} \quad \left[ \frac{N \cdot s}{m^5} \right]
\]
We now introduce the \textbf{specific acoustic impedance}, defined as:
\[
Z_s = \frac{P}{u} \quad \left[ \frac{N \cdot s}{m^3} \right]
\]
where \( u \) is the \textbf{particle velocity}.  
In this formulation, we refer to molecular velocity rather than volume velocity. The dimension of the specific acoustic impedance is the \textbf{Rayl}.
\[
Z_s = \frac{P}{u} = \frac{P}{\frac{U}{S}} = \frac{P}{U} \cdot S \Rightarrow Z_s = Z_A \, S
\]
Thus, the specific acoustic impedance is obtained by multiplying the acoustic impedance by the surface area of the radiating element.



\subsection{Signal Transformer}

In many situations, we need to convert an electrical signal from one configuration \((v, i)\) to another, while keeping the same power:
\[
P_E = v \, i
\]
A \textbf{transformer} is an electrical device designed for this purpose. It allows conversion of voltage and current levels without altering the total electrical power.

\begin{itemize}
    \item The \textbf{primary coil} (with $n_p$ coils/turns) generates a magnetic field \( B \) when a current \( i \) flows through it.
    \item The \textbf{secondary coil} (with $n_s$ coils/turns) , exposed to the magnetic flux variation \( \Delta \Phi_B \), experiences an induced voltage \( v \) and current \( i \).
\end{itemize}

When the primary coil carries current, it produces a magnetic field that couples with the secondary coil, inducing a current in it.  
The direction (or phase) of both voltages and currents is indicated by the \textbf{dots} on the circuit schematic.

\begin{figure}[H]
    \centering
    \includegraphics[width=0.55\textwidth]{signal_transformer.png}
    \caption{Transformer coupling between primary and secondary coils. The dots indicate phase correspondence.}
\end{figure}

The \textbf{transformer law} relates voltage and turns ratio:
\[
\frac{v_p}{v_s} = \frac{n_p}{n_s} = k
\]
where \( k \) is called the \textbf{turn ratio}.

\begin{itemize}
    \item If \( n_s > n_p : 
    \begin{cases}
    v_s > v_p\\
    i_s < i_p
    \end{cases}
    \quad \text{(step-up transformer)}
    \)
    \item If \( n_s < n_p : 
    \begin{cases}
    v_s < v_p\\
    i_s > i_p
    \end{cases}
    \quad \text{(step-down transformer)}
    \)
\end{itemize}
Since the transformer conserves input power:
\[
\textcolor{red}{
v_p \, i_p = v_s \, i_s
\quad \Rightarrow \quad
\frac{i_s}{i_p} = \frac{v_p}{v_s} = k
}
\]

\begin{figure}[H]
    \centering
    \begin{minipage}[c]{0.6\textwidth}
        The transformer is a \textbf{passive component}: it does not amplify power, 
        but it enables voltage and current adaptation between circuits.

        \medskip

        \medskip

        A common application is found in power distribution systems, where transformers 
        convert high transmission voltages into domestic voltage levels.
    \end{minipage}\hfill
    \begin{minipage}[c]{0.35\textwidth}
        \centering
        \includegraphics[width=\textwidth]{power_transfer_block.png}
        \caption{Typical example of a power transformer after a transmission line: 
        voltage reduction from 380~kV to 220~V.}
    \end{minipage}
\end{figure}


\subsubsection{Load Impedance Transformation}

An important aspect of the signal transformer is that it modifies the load impedance seen by the electrical power source.

We have:
\[
Z_{LS} = \frac{v_s}{i_s}
\]
We also know that:
\[
Z_p = \frac{v_p}{i_p} = \frac{k \, v_s}{i_s / k} = k^2 \frac{v_s}{i_s} 
\quad \Rightarrow \quad Z_p = k^2 Z_{LS}
\]
Thus:
\[
k > 1 \Rightarrow Z_{LS} < Z_p, \qquad k < 1 \Rightarrow Z_{LS} > Z_p
\]

This relation shows that the transformer scales the impedance by the square of the turn ratio \( k \).

\bigskip
When using a transformer, a key design question arises: given the power \( P \), are we satisfied with the actual impedance, or do we need to adjust it?


\textbf{Example:}  Consider a valve amplifier driving a loudspeaker with:
\[
P = 100\,W, \quad R_{load} = 4\,\Omega, \quad k = 20
\]
\begin{figure}[H]
    \centering
    \begin{minipage}[c]{0.55\linewidth}
We admit:
\[
i_p = 250 \times 10^{-3} \, A, \qquad v_p = 400\,V
\]
Then:
\[
v_s = \frac{v_p}{k} = 20\,V, \qquad i_s = k \, i_p = 5\,A
\]
The primary coil impedance is therefore:
\[
Z_p = k^2 Z_{LS} = 20^2 \times 4 = 1600\,\Omega = 1.6\,k\Omega
\]
    \end{minipage}\hfill
    \begin{minipage}[c]{0.4\linewidth}
        \centering
        \includegraphics[width=\linewidth]{valve_amp_to_loudspeaker.png}
    \caption{Schematic of the example.}
    \end{minipage}
\end{figure}
Considering a Thevenin source at the input of a transformer, we can describe an equivalent circuit by replacing the transformer with its reflected impedance \( Z_p = k^2 Z_{LS} \). This allows us to analyze the passive response by applying a voltage test \( v_s \) and measuring the resulting current \( i_s \).


\begin{figure}[H]
    \centering
    \includegraphics[width=0.8\textwidth]{Transformer_Thevenin_equiv.png}
    \caption{Thevenin equivalent of the signal tranformer.}
\end{figure}

We have for equivalent voltage source and impedance in the Thevenin representation: 
\[
\begin{cases}
    & v_{eq} = \frac{v_0}{k}\\
    & Z_{eq} = \frac{v_s}{i_s} = \frac{v_p/k}{i_p \cdot k} = \frac{Z_0}{k^2}
\end{cases}
\]

\subsection{Gyrator}

A \textbf{gyrator} is a two-port device that converts one \((v, i)\) configuration into another while maintaining the same electrical power.  
Unlike the transformer, the gyrator relates different electrical quantities through the \textbf{transconductance} \( g \), expressed in Siemens (\( S = \Omega^{-1} \)).
\begin{figure}[H]
    \centering
    \includegraphics[width=0.4\linewidth]{gyrator_model.png}
    \caption{Model of the gyrator.}
    \label{fig:placeholder}
\end{figure}
\textbf{Gyrator laws:} $ 
i_s = g \, v_p, \quad v_s = \dfrac{i_p}{g}$

Hence, we also observe that: $ v_p = \dfrac{i_s}{g}, \quadi_p = v_s \, g $

We know that:
\[
Z_p = \frac{v_p}{i_p} = \frac{i_s / g}{g \, v_s} = \frac{1}{g^2} \frac{i_s}{v_s} = \frac{1}{g^2} \frac{1}{Z_{LS}} 
\quad \Rightarrow \quad Z_p = \frac{1}{g^2} Y_{LS}
\]
Thus, the gyrator inverts impedance and converts it according to:
\[
Z_p = \frac{1}{g^2 Z_{LS}}
\]
Examples :
\begin{itemize}
    \item \textbf{Load is a resistor}, i.e. if \( Z_{LS} = R \), then:
\[
Z_p = \frac{1}{g^2 R} = \frac{1}{g^2} G
\]
where \( G = \frac{1}{R} \) is the conductance.
\item \textbf{Load is an inductor}, i.e. if \( Z_{LS} = j \omega L \), then:
\[
Z_p = \frac{1}{g^2 j \omega L} = \frac{1}{j \omega C_{eq}} \quad \text{with} \quad C_{eq} = g^2 L
\]
Thus, an \textbf{inductor is transformed into a capacitor} (behaviorally).
\item  \textbf{Load is a capacitor}, i.e. if \( Z_{LS} = \frac{1}{j \omega C} \), then:
\[
Z_p = \frac{1}{g^2} \, j \omega C = j \omega L_{eq} \quad \text{with} \quad L_{eq} = \frac{1}{g^2 C}
\]
Hence, a \textbf{capacitor is transformed into an inductor} (behaviorally).
\end{itemize}


\subsection{Loudspeaker Pipeline}

A loudspeaker can be modeled as a sequence of three transductions:
\[
\text{Electrical Domain} \longrightarrow \text{Mechanical Domain} \longrightarrow \text{Acoustical Domain}
\]

\begin{figure}[H]
    \centering
    \includegraphics[width=0.9\textwidth]{loudspeaker_pipeline.png}
    \caption{Schematic representation of the loudspeaker transduction chain.}
\end{figure}


\subsubsection{Electrical to Mechanical Transduction}

The conversion from electrical to mechanical energy occurs within the voice coil, where electrical current interacts with the magnetic field.  
This stage can be represented either by a \textbf{transformer} or by a \textbf{gyrator}, depending on the analogy used.

\begin{figure}[H]
    \centering
    \includegraphics[width=0.5\textwidth]{electrical_mechanical_transformer.png}
    \caption{Electrical-to-mechanical transduction block.}
\end{figure}

\noindent In the \textbf{admittance analogy}, we use a transformer: $ F_c = B \, l \, i,\quad 
u_c = \dfrac{1}{B \, l} \, v $

In the \textbf{impedance analogy}, the same relationships hold: $  F_c = B \, l \, i,\quad u_c = \dfrac{1}{B \, l} \, v $
\begin{figure}[H]
    \centering
    \includegraphics[width=0.7\linewidth]{imageE_M_analogies.png}
    \caption{Enter Caption}
    \label{fig:placeholder}
\end{figure}

The laws do not change, but the nature of the coupling does: under the impedance analogy, this corresponds to a \textbf{gyrator}.  
We therefore identify the gyrator’s transconductance as:
\[
g = \frac{1}{B \, l}
\]
where \( B \) is the magnetic flux density and \( l \) the effective length of the wire in the magnetic field.


\subsubsection{Mechanical to Acoustical Transduction}
\begin{figure}[H]
    \centering
    \includegraphics[width=0.3\linewidth]{M_A_Block.png}
    \caption{Mechanical-to-acoustical block.}
    \label{fig:placeholder}
\end{figure}
The conversion from mechanical motion to acoustic pressure occurs through the diaphragm surface.  

Using the impedance analogy: 
\[
P = \dfrac{F}{S_D},\quad U = u \, S_D
, \quad \text{where \( S_D \) is the diaphragm surface area}
\]
In the admittance analogy, the relationships combine into:
\[
P \, U = F \, u,
\]
which corresponds to a \textbf{signal transformer} model.

\begin{figure}[H]
    \centering
    \includegraphics[width=0.85\textwidth]{M_A_Analogies.png}
    \caption{Mechanical-to-acoustical transduction modeled as transformer (admittance analogy) or gyrator (impedance analogy).}
\end{figure}


In summary:
\begin{itemize}
    \item In the \textbf{admittance analogy}, the transduction behaves like a transformer.
    \item In the \textbf{impedance analogy}, the transduction behaves like a gyrator.
\end{itemize}


\subsection{Mechanical Section of the Loudspeaker}

\begin{figure}[H]
    \centering
    \begin{minipage}[c]{0.56\linewidth}
The basket of the loudspeaker can be modeled as a \textbf{hard immobile wall}.  
The diaphragm, attached to the basket by the suspension, can be represented by a \textbf{mass} \( M_{MD} \) of surface \( S_D \), moving at velocity \( u_C \) — the same as that of the voice coil.

\medskip
The suspension behaves as a \textbf{mechanical compliance} \( C_{MS} \) (elastic part) in parallel with a \textbf{mechanical resistance} \( R_{MS} \) (viscous part), accounting for losses in the suspension.

\medskip
The \textbf{mechanical conductance} of the suspension is defined as:
\[
G_{MS} = \frac{1}{R_{MS}}
\]
The \textbf{coil’s Lorentz force} drives the motion of the diaphragm:
\[
F_C = B \, l \, i
\]
The \textbf{spring} reacts with a restoring force that acts in the opposite direction:
\[
F_S = \frac{1}{C_M} \, x
\]
    \end{minipage}\hfill
    \begin{minipage}[c]{0.43\linewidth}
        \centering
        \includegraphics[width=\linewidth]{mechanical_loudspeaker_model.png}
    \caption{Mechanical representation of the loudspeaker suspension and diaphragm.}
    \end{minipage}
\end{figure}
Differentiating the spring's restoring force expression, we have:
\[
\frac{dF}{dt} = \frac{1}{C_M} \, u \quad \Leftrightarrow \quad u = C_M \, \frac{dF}{dt}
\]
The system reaches equilibrium when \( F = F_S \).

The \textbf{damper} produces a resistive force:
\[
F_D = R_M \, u \quad \Leftrightarrow \quad u = \frac{1}{R_M} \, F_D
\]
The system becomes static when \( F = F_D \).

The\textbf{ diaphragm’s mass} reacts to an applied force as:
\[
F_M = M_M \, \frac{du}{dt} \quad \Leftrightarrow \quad u = \frac{1}{M_M} \int_0^t F_M(\tau) \, d\tau
\]
According to D’Alembert’s principle, the sum of all internal and external forces must balance:
\[
F = F_S + F_D + F_M
\]
Substituting the expressions for each term:
\[
F = \frac{1}{C_M} \int_0^t u(\tau) \, d\tau + R_M \, u + M_M \, \frac{du}{dt}
\]


\subsubsection{Admittance and Impedance Analogies in the Mechanical Section}

Working in the \textbf{Admittance Analogy}: $  u_c \equiv v,\quad
F_c \equiv i $.

In this analogy:
\begin{itemize}
    \item The \textbf{mass} behaves as a \textbf{capacitor}.
    \item The \textbf{suspension} (spring and mechanical losses) behaves as an \textbf{inductor} and a \textbf{conductance}.
\end{itemize}

We obtain an \textbf{RLC circuit in parallel}, since all components share the same velocity (analogous to voltage in the electrical domain).

Working in the \textbf{Impedance Analogy}: $ u_c \equiv i,\quad F_c \equiv v $

In this analogy:
\begin{itemize}
    \item The \textbf{mass} behaves as an \textbf{inductor}.
    \item The \textbf{suspension} (spring and mechanical resistance) behaves as a \textbf{capacitor} and a \textbf{resistor}.
\end{itemize}

We obtain an \textbf{RLC circuit in series}, since all components share the same speed (analogous to current in the electrical domain).

\begin{figure}[H]
    \centering
    \includegraphics[width=0.7\textwidth]{mechanical_section_analogies.png}
    \caption{Admittance analogy (left): mechanical section modeled as an RLC parallel circuit. Impedance analogy (right): mechanical section modeled as an RLC series circuit.}
\end{figure}

\clearpage

\textbf{Example of the Entire Signal Path: }

\begin{figure}[H]
    \centering
    \includegraphics[width=0.95\textwidth]{loudspeaker_full_model.png}
    \caption{Complete equivalent model of the loudspeaker: electrical, mechanical, and acoustical domains.}
\end{figure}


The voice coil is modeled by a RL series circuit.  
The \textbf{electrical-to-mechanical} transduction is represented by a \textbf{signal transformer}, while the \textbf{mechanical section} of the loudspeaker is modeled using the \textbf{admittance analogy} (RLC in parallel).  
Finally, the \textbf{mechanical-to-acoustical} transduction is also modeled by a signal transformer.
\[
Z_{AR} = \frac{\tilde{P}}{\tilde{U}}
\]
is the \textbf{radiation acoustic impedance}, which relates the sound pressure at the diaphragm to the corresponding volume velocity.

\subsubsection{Lumped Elements Model}

In many systems (such as transmission lines), it is useful to approximate distributed physical phenomena with equivalent lumped elements.  
This approach replaces a spatially distributed system by discrete resistors, inductors, and capacitors that reproduce its behavior.

\begin{itemize}
    \item A single RLC circuit can provide a first-order approximation of the system’s behavior.
    \item Each RLC element represents the resistive, inductive, and capacitive properties of a small section of the system.
    \item The complete model is obtained by connecting multiple RLC sections in series.
\end{itemize}

\begin{figure}[H]
    \centering
    \includegraphics[width=0.75\textwidth]{lumped_elements_model.png}\caption{Transition from a distributed system to a lumped-element RLC approximation.}
\end{figure}

For high frequencies (when the wavelength \( \lambda \) becomes comparable to the system’s dimensions), a finer subdivision of the wire into many RLC portions is required.  
In this case, distributed effects dominate, and the lumped model must be extended.

\begin{figure}[H]
    \centering
    \includegraphics[width=0.7\linewidth]{lumped_elements_high_freq.png}
    \caption{Long wire line divided in many many lumped elements (RLC sections).}
    \label{fig:placeholder}
\end{figure}

\subsubsection{Pressure Waves}

When an electrical signal excites the coil and moves the diaphragm, it creates a pressure increase at the front of the loudspeaker and a pressure decrease at the back — the two being in phase opposition.

\begin{figure}[H]
    \centering
    \begin{minipage}[c]{0.25\textwidth}
        \centering
        \includegraphics[width=\textwidth]{loudspeaker_pressure_front_back.png}
        \caption{Front and back acoustic pressures generated by the diaphragm (in phase opposition).}
    \end{minipage}
    \hfill
    \begin{minipage}[c]{0.7\textwidth}
        \centering
        \includegraphics[width=\textwidth]{impedance_redraw.png}
        \caption{Alternate modeling of the front and back air impedances.}
    \end{minipage}
\end{figure}


As the front and back environments are different, their corresponding air impedances \( Z_{AF} \) and \( Z_{AB} \) are not identical.

The equivalent acoustic impedance seen by the diaphragm is:
\[
Z_{AR} = Z_{AB} + Z_{AF}
\]
If we assume symmetry between front and back (\( Z_{AB} = Z_{AF} \)), we can simplify:
\[
Z_{AR} = 2Z_{AF} \quad \text{or equivalently} \quad Y_{AR} = \frac{Y_{AF}}{2}
\]


\subsection{From acoustical to mechanical domain}

Starting from the previous observations, we can define the remaining part of the model. We have:
\[
Y_{MR} = \frac{u}{F} = \frac{U/S_D}{P \cdot S_D} = \frac{1}{S_D^2} \frac{U}{P} 
\quad \Rightarrow \quad  Y_{MR} = \frac{1}{S_D^2} Y_{AR}
\]
Note that:
\[
Z_{MR} = \frac{F}{u} = \frac{1}{Y_{MR}} = S_D^2 \, Z_{AR}
\]
\textbf{NB:} MR stands for \textit{mechanical radiation}.

\begin{figure}[H]
    \centering
    \includegraphics[width=0.6\textwidth]{mechanical_radiation_model.png}
    \caption{Model of the acoustic radiation impedance of the loudspeaker.}
\end{figure}

\subsection{Mechanical model}

The loudspeaker is enclosed in a box, and the air inside behaves as a spring (compliance). The equivalent mechanical circuit can therefore be modeled as a mass-spring-damper system where the diaphragm motion is characterized by the mechanical admittance \( Y_{MR} \).

\begin{figure}[H]
    \centering
    \includegraphics[width=0.45\textwidth]{mechanical_model.png}
    \caption{Mechanical model adding the air inside the loudspeaker's box.}
\end{figure}


\subsubsection{From mechanical to acoustical domain}

We can now express the electrical radiation impedance as:
\[
Z_{ER} = \frac{v}{i} = \frac{u \cdot B \cdot l}{F / (B \cdot l)}
\]
Thus:
\[
Z_{ER} = (B \, l)^2 \frac{u}{F} = (B \, l)^2 Y_{MR} = (B \, l)^2 Y_{MR} = \left( \frac{B \, l}{S_D} \right)^2 Y_{AR}
\]
\textbf{NB:} ER stands for \textit{electrical radiation}.

\begin{figure}[H]
    \centering
    \includegraphics[width=0.6\textwidth]{electrical_radiation_model.png}
    \caption{Model of the electrical radiation impedance of the loudspeaker.}
\end{figure}


\subsubsection{Mass of the diaphragm and voice coil}

The mechanical equation governing the diaphragm and voice coil is derived from Newton’s second law:
\[
F = M_{MD} \frac{du}{dt} \quad \Rightarrow \quad i = C \frac{dv}{dt}
\]
In the sinusoidal domain:
\[
F = j \omega M_{MD} u \quad \text{and} \quad i = j \omega C \, v
\]
Using the electromechanical analogy, we have:
\[
\frac{u}{F} = \frac{1}{j \omega M_{MD}}, \qquad \frac{v}{i} = (B l)^2 \frac{u}{F} = (B l)^2 \frac{1}{j \omega M_{MD}}
\qquad \Rightarrow \quad 
\frac{v}{i} = \frac{1}{j \omega} \frac{M_{MD}}{(B l)^2}
\]
This corresponds to a capacitor of capacitance:
\[
C_{ED} = \frac{M_{MD}}{(B l)^2}
\]
Hence, a smaller diaphragm mass \( M_{MD} \) yields a smaller equivalent electrical capacitance \( C_{ED} \). Since we usually aim for \( C_{ED} \) in the order of nanofarads, an ideal loudspeaker should have a small moving mass.

\begin{figure}[H]
    \centering
    \begin{minipage}[c]{0.45\textwidth}
        \centering
        \includegraphics[width=\textwidth]{Diaphragm_capacitor_M_MD.png}
    \end{minipage}
    \hspace{2cm}
    \begin{minipage}[c]{0.2\textwidth}
        \centering
        \includegraphics[width=\textwidth]{Diaphragm_capacitor_C_ED.png}
    \end{minipage}
    \caption{Electrical Diaphragm capacitance in the acoustic domain (left) and electric domain (right).}
\end{figure}

\subsubsection{Suspension compliance}

We denote \( C_{MS} \) as the mechanical compliance of the loudspeaker suspension system.  
From the fundamental definition of compliance, we have:
\[
x = C_{MS} \cdot F
\quad \Rightarrow \quad u = C_{MS} \frac{dF}{dt}
\equiv v = L \frac{di}{dt}
\]
using the electromechanical admittance analogy.

Assuming a sinusoidal steady-state regime, we obtain:
\[
\frac{u}{F} = j \omega C_{MS}, \quad \text{since \(\dfrac{dF}{dt} = j \omega F\).}
\]
Therefore, using the transformer law linking mechanical and electrical domains, we can express:
\[
\frac{v}{i} = (B l)^2 \frac{u}{F}
\Rightarrow \frac{v}{i} = j \omega (B l)^2 C_{MS}
\]
which corresponds to an inductor:
\[
L_{ES} = (B l)^2 C_{MS}
\]
\textbf{NB:} This inductance represents the electrical equivalent of the suspension compliance.

\begin{figure}[H]
    \centering
    \begin{minipage}[c]{0.5\textwidth}
        \centering
        \includegraphics[width=\textwidth]{electrical_suspension_inductance.png}
    \end{minipage}
    \hfill
    \begin{minipage}[c]{0.3\textwidth}
        \centering
        \includegraphics[width=\textwidth]{electrical_suspension_inductance_2.png}
    \end{minipage}
    \caption{Electrical suspension inductance (left) and resitance (right).}
\end{figure}


\subsubsection{Suspension conductance}

The mechanical losses of the suspension can be modeled as a conductance \( G_{MS} \) defined by:
\[
G_{MS} = \frac{u}{F}
\]
Using the transformer analogy again, we have:
\[
\frac{v}{i} = (B l)^2 \frac{u}{F}
\Rightarrow \frac{v}{i} = (B l)^2 G_{MS}
\]
Therefore, the suspension introduces an electrical resistance:
\[
R_{ES} = \frac{1}{(B l)^2 G_{MS}} = \frac{(B l)^2}{R_{MS}}
\]
where \( R_{MS} = \frac{1}{G_{MS}} \) is the mechanical resistance.


\begin{figure}[H]
    \centering
    \begin{minipage}[c]{0.5\textwidth}
        \centering
        \includegraphics[width=\textwidth]{electrical_suspension_resistance_0.png}
    \end{minipage}
    \hfill
    \begin{minipage}[c]{0.3\textwidth}
        \centering
        \includegraphics[width=\textwidth]{electrical_suspension_resistance.png}
    \end{minipage}
    \caption{Electrical suspension inductance (left) and resitance (right).}
\end{figure}


The higher the electrical resistance \( R_{ES} \), the lower the mechanical resistance \( R_{MS} \), and therefore the greater the mechanical damping of the loudspeaker system.


\subsubsection{Overall equivalent model}

By combining all the previous elements, the mechanical representation of the loudspeaker system becomes:

\begin{figure}[H]
    \centering
    \includegraphics[width=0.9\textwidth]{overall_mechanical_model.png}
    \caption{Circuit of the full loudspeaker's model.}
\end{figure}

To simplify the representation, we introduce two equivalent impedances:
\[
Z_{ES} = R_{ES} + Z_{L_{ES}}
\]
which corresponds to the \textbf{suspension-related impedance}, and
\[
Z_{EM} = Z_{L_{E_C}} \,//\, R_{ES} \,//\, Z_{C_{ED}} \,//\, Z_{ER}
\]
which corresponds to the \textbf{mass-related impedance}.  
The total equivalent impedance of the electro-mechanical system is therefore:
\[
Z_E = Z_{ES} + Z_{EM}
\]
\begin{figure}[H]
    \centering
    \includegraphics[width=0.7\textwidth]{overall_model_2.png}
    \caption{Circuit of the model with equivalent impedances.}
\end{figure}


We observe that the electrical analogy of the loudspeaker system can be expressed as a parallel \textbf{RLC circuit}, which exhibits a resonance at a specific frequency \( f_s \).  
This resonance corresponds to the natural oscillation frequency of the moving mass and suspension system.

\subsection{Quick summary about basics}

For reference, consider a general RLC circuit in parallel configuration:

The impedance of a parallel RLC circuit is given by:
\[
Z_{RLC} = \frac{1}{\frac{1}{R} + j\omega C - \frac{j}{\omega L}}
\]
The magnitude of this impedance reaches a maximum at the \textbf{resonant frequency}:
\[
\omega_s = \frac{1}{\sqrt{L C}} \quad \Rightarrow \quad f_s = \frac{1}{2\pi\sqrt{L C}}
\]
At this frequency, the inductive and capacitive reactances cancel each other (\( \omega L = 1 / (\omega C) \)), and the total impedance becomes purely resistive.

Below resonance (\( \omega < \omega_s \)), the circuit behaves predominantly capacitive, while above resonance (\( \omega > \omega_s \)), it behaves inductively.  
This same principle applies to the loudspeaker’s electro-mechanical equivalent circuit.

\subsection{Summary: LC series and parallel circuits}

We consider a series and a parallel LC circuit.  
At resonance, the magnitudes of the inductive and capacitive impedances are equal:
\[
|Z_L| = |Z_C| \quad \Rightarrow \quad j \omega_s L = \frac{1}{j \omega_s C} \quad \Rightarrow \quad \omega_s = \frac{1}{\sqrt{L C}}
\]

\subsubsection{LC Series Circuit}

For the series configuration:
\[
Z_S = \frac{v}{i} = Z_L + Z_C = j \omega L + \frac{1}{j \omega C} = j \left( \omega L - \frac{1}{\omega C} \right)
\]
Rewriting:
\[
Z_S = j \frac{\omega^2 L C - 1}{\omega C} = j \frac{1 - \omega_s^2 / \omega^2}{\omega C}
\]
At the resonance frequency \( \omega = \omega_s \), we have:
\[
Z_S(\omega_s) = 0 \quad \Rightarrow \quad \text{short circuit (minimum impedance).}
\]
The magnitude of the impedance is given by:
\[
|Z_S| = \frac{|1 - (\omega / \omega_s)^2|}{\omega C}
\]
At low frequencies, the capacitor dominates (\(Z_S \approx 1 / j\omega C\)), while at high frequencies, the inductor dominates (\(Z_S \approx j\omega L\)).

\subsubsection{LC Parallel Circuit}

For the parallel configuration:
\[
Z_P = \frac{v}{i} = \frac{Z_L Z_C}{Z_L + Z_C}
\]
Substituting \( Z_L = j \omega L \) and \( Z_C = \frac{1}{j \omega C} \):
\[
Z_P = \frac{j \omega L \cdot \frac{1}{j \omega C}}{j \omega L + \frac{1}{j \omega C}} = \frac{L / C}{j (\omega^2 L C - 1)} = - j \frac{\omega L}{1 - \omega^2 L C}
\]
At the resonance frequency \( \omega = \omega_s \):
\[
Z_P(\omega_s) \rightarrow \infty \quad \Rightarrow \quad \text{open circuit (maximum impedance).}
\]
The magnitude of the impedance is:
\[
|Z_P| = \frac{\omega L}{|1 - (\omega / \omega_s)^2|}
\]


\subsubsection{Interpretation}

\begin{itemize}
    \item In a \textbf{series LC circuit}, the impedance reaches a minimum (short circuit) at resonance.  
    \item In a \textbf{parallel LC circuit}, the impedance reaches a maximum (open circuit) at resonance.  
    \item The resonance frequency is the same for both configurations: $ f_s = \frac{1}{2 \pi \sqrt{L C}} $
    \item The behavior around resonance explains the presence of peaks or dips in the loudspeaker’s impedance curve, which corresponds to the mechanical resonance of the diaphragm–suspension system.
\end{itemize}


\subsection{RLC Parallel Configuration and Frequency Response}

\begin{figure}[H]
    \centering
    \begin{minipage}[c]{0.6\textwidth}
        The impedance of a parallel RLC circuit can be expressed as:
\[
Z_P = \frac{v}{i} = \left( Y_L + Y_C + Y_R \right)^{-1} \Rightarrow Z_P = \left( \frac{1}{j \omega L} + j \omega C + \frac{1}{R} \right)^{-1}
\]
Simplifying:
\[
Z_P = \left( \frac{1 + j \omega^2 L C + \frac{j \omega L}{R}}{j \omega L} \right)^{-1}
\Rightarrow Z_P = \frac{j \omega L R}{R(1 - \omega^2 L C) + j \omega L}
\]
    \end{minipage}
    \hfill
    \begin{minipage}[c]{0.35\textwidth}
        \centering
        \includegraphics[width=\textwidth]{RLC parallel.png}
    \caption{RLC parralel circuit.}
    \end{minipage}
\end{figure}
Rewriting:
\[
Z_P = \frac{j \omega L R}{R(1 - \omega^2 L C) + j \omega L} = \frac{j \omega L R}{R - \omega^2 R L C + j \omega L}
\Rightarrow Z_P = \frac{j \omega L R}{R(1 - \omega^2 L C + j \omega L / R)}
\]
Thus, the general impedance expression becomes:
\[
Z_P = \frac{j \omega L}{(1 - \omega^2 L C) + j \omega L / R}, \quad 
\omega_s = \frac{1}{\sqrt{L C}}
\]
Substituting this into the previous expression:
\[
Z_P = \frac{j \omega L}{\left(1 - \frac{\omega^2}{\omega_s^2}\right) + j \frac{\omega L}{R}}
\]
The \textbf{quality factor \(Q\) }expresses the sharpness of the resonance:
\[
\frac{1}{Q} = \frac{\omega_s L}{R} = \frac{2 \pi f_s L}{R}
\quad \Rightarrow \quad Q = \frac{R}{L \omega_s} = \frac{R}{L} \sqrt{L C} 
\quad \Rightarrow \quad  Q = \frac{R}{L \omega_s} = R \sqrt{\frac{C}{L}}
\]
which can also be expressed as:
\[
Q = \frac{f_s}{\text{BW}}
\]
where $BW$ denotes the \textbf{bandwidth} between the frequencies for which the impedance magnitude decreases by \( \frac{1}{\sqrt{2}} \) of its maximum.

\begin{figure}[H]
    \centering
    \begin{minipage}[c]{0.6\textwidth}
        Substituting \( Q \) into the impedance expression:
\[
Z_P = R \, \frac{j \frac{\omega}{\omega_s} \frac{1}{Q}}{1 - \left( \frac{\omega}{\omega_s} \right)^2 + j \frac{\omega}{\omega_s} \frac{1}{Q}}
\]
At resonance (\( \omega = \omega_s \)):
\[
Z_P(\omega_s) = R
\]
The impedance magnitude reaches its maximum, indicating the \textbf{resonant frequency} of the system.  
    \end{minipage}
    \hfill
    \begin{minipage}[c]{0.35\textwidth}
        \centering
        \includegraphics[width=\textwidth]{Q_Factor_Graph.png}
    \caption{Influence of $Q$ on the magnitude response.}
    \end{minipage}
\end{figure}

Thus, the \textbf{electrical motion impedance} of the loudspeaker’s equivalent circuit is:
\[
Z_{EM} = R \, \frac{j \frac{\omega}{\omega_s} \frac{1}{Q}}{1 - \left( \frac{\omega}{\omega_s} \right)^2 + j \frac{\omega}{\omega_s} \frac{1}{Q}}
\]


The electrical motion impedance \( Z_{EM} \) is defined as:
\[
Z_{EM} = \left( j \omega C_{ED} + \frac{1}{j \omega L_{EC}} + \frac{1}{R_{ES}} + \frac{1}{Z_{ER}} \right)^{-1}
\]
and we recall that:
\[
\frac{1}{Z_{ER}} = j \omega \frac{2 M_{M1}}{(B l)^2}
\]
We set the total moving mass as:
\[
M_{MS} = M_{MD} + 2 M_{M1}
\quad \Rightarrow \quad 
Z_{EM} = \left( j \omega \frac{M_{MS}}{(B l)^2} + (B l)^2 C_{MS} + \frac{1}{(B l)^2 G_{MS}} \right)^{-1}
\]

\begin{figure}[H]
    \centering
    \begin{minipage}[c]{0.45\textwidth}
        After simplification, we obtain:
\[
Z_{EM} = \frac{(B l)^2}{R_{MS}} \cdot \frac{j \dfrac{1}{Q_{MS}} \dfrac{\omega}{\omega_s}}{1 - \left( \dfrac{\omega}{\omega_s} \right)^2 + j \dfrac{1}{Q_{MS}} \dfrac{\omega}{\omega_s}}
\]
where:
\[
\omega_s = \frac{1}{\sqrt{M_{MS} C_{MS}}}, \qquad Q_{MS} = \frac{1}{R_{MS}} \sqrt{\frac{M_{MS}}{C_{MS}}}
\]
    \end{minipage}
    \hfill
    \begin{minipage}[c]{0.5\textwidth}
        \centering
        \includegraphics[width=\textwidth]{Z_EM(f).png}
    \caption{Graph of $|Z_{EM}(f)$ and it's different elements.}
    \end{minipage}
\end{figure}


The \textbf{total electrical impedance of the loudspeaker} can be expressed as:
\[
Z_E = Z_{ES} + Z_{EM} = R_{EC} + j \omega L_{EC} + Z_{EM}
\]
where:
\begin{itemize}
    \item \( R_{EC} \) is constant with frequency (electrical losses),
    \item \( j \omega L_{EC} \) increases linearly with frequency (voice coil inductance),
    \item \( Z_{EM} \) represents the mechanical resonance of the diaphragm-suspension system, behaving as a damped mass-spring impedance.
\end{itemize}

At this stage, \textbf{we can identify the main physical quantities describing the loudspeaker}:
\begin{itemize}
    \item \textbf{Acoustic power and efficiency:} $  \varepsilon = \frac{P_A}{P_E} = \frac{p \cdot U}{v \cdot i} $ 
    \item \textbf{Sound Pressure Level (SPL)}
    \item \textbf{Frequency response}
\end{itemize}

\subsection{From electrical to acoustical domain}

In a loudspeaker, the current flowing through the voice coil generates a magnetic field, which interacts with the magnet and drives the diaphragm.  
The electrical power \( P_E \) is thus converted into mechanical power \( P_M \), which is in turn radiated as acoustic power \( P_A \), corresponding to the movement of air particles.
\[
P_E = v(t) \cdot i(t) \quad \Rightarrow \quad P_M = F(t) \cdot u(t) \quad \Rightarrow \quad P_A = p(t) \cdot U(t)
\]
To describe an audio signal in terms of power, we define an equivalent constant power \( \bar{P} \) that represents the same total energy as the time-varying power \( P(t) \).

The energy of the variable and constant signals are:
\[
E_V = \int_0^T P(t)\,dt, \qquad E_C = \bar{P} \cdot T
\]
By equating \( E_V = E_C \), we get:
\[
\bar{P} = \frac{1}{T} \int_0^T P(t)\,dt
\]
which corresponds to the time-averaged power.

The instantaneous electrical power associated with a resistor is:
\[
P(t) = v(t) \cdot i(t) = \frac{v^2(t)}{R}
\]
Hence, the mean power becomes:
\[
\bar{P} = \frac{1}{R} \cdot \frac{1}{T} \int_0^T v^2(t)\,dt = \frac{\bar{v^2}}{R} = \frac{v_{RMS}^2}{R}
\]
since \( v_{RMS} = \sqrt{\bar{v^2}} \), corresponding to the equivalent constant voltage that produces the same power transfer as \( v(t) \).

In a sinusoidal regime:
\[
v(t) = V_P \sin(2\pi f t + \phi) \quad \Rightarrow \quad v_{RMS} = \frac{V_P}{\sqrt{2}} \quad \Rightarrow \quad v_{RMS|dB} = V_P - 3\,\mathrm{dB}
\]
Thus:
\[
\bar{P} = \frac{v_{RMS}^2}{R} = \frac{V_P^2}{2R}
\]
and the peak power is:
\[
P_{peak} = \frac{V_P^2}{R} = 2 \bar{P}
\]

Note that real musical signals are more complex than a single sinusoid, meaning that both peak and RMS values (for voltage or power) vary over time.  
To describe this difference, we define the \textbf{crest factor (CF)} as the ratio between peak power and mean power:
\[
CF = \frac{P_{peak}}{\bar{P}}
\]
Typical values:
\begin{cases}
    & CF = 2 \; (+3\,\mathrm{dB}) \quad \text{for sinusoidal signals} \\
    & CF \geq 4 \; (\geq +6\,\mathrm{dB}) \quad \text{for generic audio signals}
\end{cases}

When expressed in decibels:
\[
CF_{dB} = 10 \log_{10}(CF)
\]

\subsubsection{Power and impedance in the electrical domain}

We consider a sinusoidal steady-state regime:
\[
v(t) = V \sin(\omega t), \qquad i(t) = I \sin(\omega t + \beta)
\]
where \( \beta \) is the phase difference between voltage and current.

The instantaneous power is given by:
\[
W_E(t) = v(t) \cdot i(t) = V I \sin(\omega t) \sin(\omega t + \beta)
\]
Using trigonometric identities:
\[
W_E(t) = V I \left[ \frac{1}{2} \cos(\beta) - \frac{1}{2} \cos(2 \omega t + \beta) \right]
\quad \Rightarrow \quad  W_E(t) = \frac{V I}{2} \cos(\beta) - \frac{V I}{2} \cos(2 \omega t + \beta)
\]
Expressing in RMS quantities:
\[
W_E(t) = V_{RMS} I_{RMS} \cos(\beta) - V_{RMS} I_{RMS} \cos(2 \omega t + \beta)
\]
\begin{figure}[H]
    \centering
    \begin{minipage}[c]{0.6\textwidth}    
The first term \( V_{RMS} I_{RMS} \cos(\beta) \) represents the \textbf{average power} \( \bar{P_E} \), constant in time,  
while the second term represents the \textbf{oscillating component} around the mean value.
\[
\bar{P_E} = V_{RMS} I_{RMS} \cos(\beta)
\]
The maximum power occurs when the voltage and current are in phase (\( \beta = 0 \)):
\[
P_{E,\text{max}} = V_{RMS} I_{RMS} \left[ \cos(\beta) + 1 \right]
\]
In this condition, the system delivers maximum energy transfer from source to load.
    \end{minipage}
    \hfill
    \begin{minipage}[c]{0.35\textwidth}
        \centering
        \includegraphics[width=\textwidth]{electrical_power_wf.png}
    \caption{Instantaneous electrical power \( W_E(t) \) showing the average component and oscillating term.}
    \end{minipage}
\end{figure}

In the particular case where:
\[
\beta = \frac{\pi}{2} \quad \Rightarrow \quad \cos(\beta) = 0
\quad \Rightarrow \quad
W_E(t) = -V_{RMS} I_{RMS} \sin(2 \omega t)
\]
In this situation, there is no net transfer of power over one period (the mean power is zero).  
The system is therefore conservative — the energy oscillates between electrical and magnetic forms without being dissipated.  
For a loudspeaker, this corresponds to oscillations around the rest position, without radiating acoustic power.


\subsubsection{Complex representation of power}

For capacitors, we have:
\[
Z_C = \frac{v}{i} = \frac{1}{j \omega C} = -\frac{j}{\omega C}
\quad \Rightarrow \quad i = j \omega C \, v
\]
In such a case, the current leads the voltage by \( \beta = \frac{\pi}{2} \).

For inductors, we have:
\[
Z_L = \frac{v}{i} = j \omega L
\quad \Rightarrow \quad i = -\frac{j}{\omega L} \, v
\]
In this case, the current lags the voltage by \( \beta = -\frac{\pi}{2} \).

Using the complex representation, we can write:
\[
i(t) = i_p(t) + j \, i_q(t) = I \sin(\omega t + \beta)
\]
Hence:
\[
i(t) = I \cos(\beta) \sin(\omega t) + I \sin(\beta) \cos(\omega t)
\Rightarrow i(t) = I_p \sin(\omega t) + I_q \sin\left(\omega t - \frac{\pi}{2}\right)
\]
where \( I_p = I \cos(\beta) \) is the in-phase component and \( I_q = I \sin(\beta) \) the quadrature component.

We can now compute the instantaneous power:
\[
W(t) = v(t) \cdot i(t) = V \sin(\omega t) \left[ I_p \sin(\omega t) + I_q \sin\left(\omega t - \frac{\pi}{2}\right) \right]
\]
Expanding:
\[
W(t) = V I_p \sin^2(\omega t) + V I_q \sin\left(\omega t - \frac{\pi}{2}\right) \sin(\omega t)
\]
Using trigonometric identities:
\[
W(t) = V I_p \frac{1 - \cos(2 \omega t)}{2} + V I_q \frac{1}{2} \sin(2 \omega t)
\]
Expressed in RMS quantities:
\[
W(t) = V_{RMS} I_{p,RMS} [1 - \cos(2 \omega t)] + V_{RMS} I_{q,RMS} \sin(2 \omega t)
\]
The first term corresponds to the \textbf{active (real) power}:
\[
W_p(t) = V_{RMS} I_{p,RMS} [1 - \cos(2 \omega t)]
\]
which is in phase with the voltage and has an average value:
\[
\bar{W}_p = V_{RMS} I_{p,RMS}
\]
The second term corresponds to the \textbf{reactive (imaginary) power}:
\[
W_q(t) = V_{RMS} I_{q,RMS} \sin(2 \omega t)
\]
which is in quadrature and has zero mean value.

\subsubsection{Example: Conservative Mechanical System}

Let us consider a simple conservative mechanical system composed of a mass and a spring:

\begin{enumerate}
    \item An impulsive force moves the mass \( M \).
    \item The spring’s potential energy makes \( M \) move back.
    \item The kinetic energy of \( M \) compresses the spring.
    \item The spring energy makes \( M \) stop and move forward again.
    \item The process repeats cyclically.
\end{enumerate}

\begin{figure}[H]
\centering
\begin{minipage}[c]{0.7\textwidth}
    \centering
    \includegraphics[width=\textwidth]{mechanical_conservative_system.png}
    \caption{Equivalent mechanical model with mass \( M \) and spring representing the conservative system.}
\end{minipage}\hfill
\begin{minipage}[c]{0.25\textwidth}
    \centering
    \includegraphics[width=\textwidth]{example_equiv_circuit.png}
    \caption{Electrical analog representation with complex impedance \( Z_E \).}
\end{minipage}
\end{figure}
We have:
\[
i(t) = \frac{v(t)}{Z_E} = \frac{v(t)}{Z_E} \cdot \frac{Z_E^*}{Z_E^*} = \frac{v(t) Z_E^*}{|Z_E|^2}
\]
and since:
\[
Z_E = \mathrm{Re}[Z_E] + j\,\mathrm{Im}[Z_E]
\]
we can expand:
\[
i(t) = v(t) \cdot \frac{\mathrm{Re}[Z_E] - j\,\mathrm{Im}[Z_E]}{(\mathrm{Re}[Z_E])^2 + (\mathrm{Im}[Z_E])^2}
\]
Thus:
\[
i(t) = \frac{\mathrm{Re}[Z_E]}{|Z_E|^2} v(t) - j \frac{\mathrm{Im}[Z_E]}{|Z_E|^2} v(t)
\quad \Rightarrow \quad i(t) = i_p(t) + j \, i_q(t)
\]
where \( i_p(t) \) is the in-phase component and \( i_q(t) \) is the quadrature component of the current:
\[
i_p(t) = \frac{\mathrm{Re}[Z_E]}{|Z_E|^2} v(t), \qquad i_q(t) = \frac{\mathrm{Im}[Z_E]}{|Z_E|^2} v(t)
\]
The instantaneous power is:
\[
W_E(t) = v(t) \, i(t) = \frac{v^2(t)}{Z_E} = \frac{\mathrm{Re}[Z_E]}{|Z_E|^2} v^2(t) - j \frac{\mathrm{Im}[Z_E]}{|Z_E|^2} v^2(t)
\]
Alternatively:
\[
W_E(t) = v(t)\,i_p(t) + j\,v(t)\,i_q(t) = W_{E,p}(t) + j\,W_{E,q}(t)
\]
The average power is then:
\[
\bar{W}_E = \bar{W}_{E,p} + j\,\bar{W}_{E,q} = \bar{W}_{E,p} \quad \Rightarrow \quad \bar{W}_E = \frac{\mathrm{Re}[Z_E]}{|Z_E|^2} \bar{v^2}
\]
Expressed in RMS terms:
\[
\bar{W} = \frac{\mathrm{Re}[Z_E]}{|Z_E|^2} V_{RMS}^2
\]
and since:
\[
I_{RMS} = \frac{V_{RMS}}{|Z_E|}
\]
we obtain:
\[
\boxed{\bar{W} = \mathrm{Re}[Z_E] \, I_{RMS}^2}
\]

\subsubsection{Power and impedance, mechanical domain}


\begin{figure}[H]
\centering
\begin{minipage}[c]{0.7\textwidth}
We use the admittance analogy
\[
F \leftrightarrow i, \qquad u \leftrightarrow v, \qquad u(t) = \frac{F(t)}{Z_M}, \quad Z_M = R_M + jX_M
\]
Decomposing \(u(t)\) into in-phase and quadrature parts:
\begin{align*}
    & u(t)=\frac{F(t)}{Z_M} = \frac{Z_M^{*}}{|Z_M|^2}F(t)\\
    \Rightarrow \quad & u(t)     
     = \frac{R_M}{|Z_M|^2}F(t) - j\,\frac{X_M}{|Z_M|^2}F(t)
     = u_p(t) + j\,u_q(t),
\end{align*}
\end{minipage}\hfill
\begin{minipage}[c]{0.25\textwidth}
\centering
\includegraphics[width=\textwidth]{mechanical_impedance_block.png}
\caption{Mechanical one-port.}
\end{minipage}
\end{figure}

Where
\[
u_p(t)=\frac{R_M}{|Z_M|^2}F(t), \qquad
u_q(t)=\frac{X_M}{|Z_M|^2}F(t).
\]
Instantaneous mechanical power:
\[
W_M(t)=F(t)\,u(t)
      = \frac{R_M}{|Z_M|^2}F^2(t)
        - j\,\frac{X_M}{|Z_M|^2}F^2(t)
        \quad \Rightarrow \quad         
      W_M(t)= R_M u^2(t) - jX_M u^2(t)
            = W_{M,p}(t)+j\,W_{M,q}(t).
\]
Average (real) power:
\[
\bar W_M=\bar W_{M,p} \ \text{since }  \overline{W}_{M,q}(t)=0
\quad \Rightarrow \quad 
\bar W_M =\frac{R_M}{|Z_M|^2}\,\overline{F^2(t)}
= R_M\,u_{\mathrm{RMS}}^{\,2}, \text{ as we have } u_{\mathrm{RMS}} = \dfrac{F_{\mathrm{RMS}}}{|Z_M|}
\]


\subsubsection{Power and impedance, acoustical domain}

\begin{figure}[H]
\centering
\begin{minipage}[c]{0.56\textwidth}
We adopt the impedance analogy:
\[
P \leftrightarrow v,\qquad U \leftrightarrow i,\qquad U(t)=\frac{P(t)}{Z_A}, \quad Z_A=R_A+jX_A
\]
Decomposition of the volume velocity:
\begin{align*}
    & U(t)=\frac{P(t)}{Z_A}
     =\frac{Z_A^{*}}{|Z_A|^2}P(t)\\
    \Rightarrow \quad & U(t) =\frac{R_A}{|Z_A|^2}P(t)-j\,\frac{X_A}{|Z_A|^2}P(t)
     =U_p(t)+j\,U_q(t)
\end{align*}
\textbf{NB:} Here \(P(t)\) denotes \emph{acoustic pressure}.
\end{minipage}\hfill
\begin{minipage}[c]{0.38\textwidth}
\centering
\includegraphics[width=\textwidth]{acoustic_port_block.png}
\caption{Acoustic one-port with impedance \(Z_A=R_A+jX_A\).}
\end{minipage}
\end{figure}

Instantaneous acoustic power:
\[
W_A(t)=P(t)\,U(t)
=\frac{R_A}{|Z_A|^2}P^2(t)-j\,\frac{X_A}{|Z_A|^2}P^2(t) 
\quad \Rightarrow \quad W_A(t) = W_{A,p}(t)+j\,W_{A,q}(t)
\]
Average power (reactive term averages to zero):
\[
\bar W_A=\bar W_{A,p}
=\frac{R_A}{|Z_A|^2}\,\overline{P^2(t)}
=\frac{R_A}{|Z_A|^2}P_{\mathrm{RMS}}^{\,2}
\quad \Rightarrow \quad
\overline{W}_A = R_A\,U_{\mathrm{RMS}}^{\,2},
\ \text{ since } U_{\mathrm{RMS}}=\dfrac{P_{\mathrm{RMS}}}{|Z_A|}
\]

\subsubsection{Specific acoustic impedance and units}

The acoustic radiation impedance depends on geometry and boundary conditions.  
The \emph{specific} acoustic impedance (per unit area) is
\[
Z_S=\frac{P}{u}=\frac{P}{U/S}=Z_A\cdot S .
\]
Units:
\[
Z_A=\frac{P}{U}\;[\,\mathrm{Pa}\cdot \mathrm{s}\,\mathrm{m}^{-3}\,]
=\frac{N/m^2}{m^3/s}=\frac{N\cdot s}{m^{5}},
\qquad
Z_S=Z_A S=\frac{\mathrm{Pa}\cdot \mathrm{s}}{m} \;\;\text{(Rayls)}.
\]

\subsubsection{Pressure–velocity relation study}

\begin{figure}[H]
\centering
\begin{minipage}[c]{0.58\textwidth}
\small
Consider an infinitesimal air element of volume \(dV=dx\,dy\,dz\) and face area \(dS=dy\,dz\).  
Pressure on the left face: \(f_1=p(x)\,dS\).  
Pressure on the right face:
\[
f_2=p(x+dx)\,dS=\Big(p(x)+\frac{\partial p}{\partial x}dx\Big)dS.
\]
Resultant force along \(x\):
\[
f=f_1-f_2=-\frac{\partial p}{\partial x}\,dx\,dS=-\frac{\partial p}{\partial x}\,dV.
\]
Newton’s law with mass \(m=\rho\,dV\) gives:
\[
f=m\,\frac{\partial \underline{u}}{\partial t}\quad\Rightarrow\quad
-\frac{\partial p}{\partial x}=\rho\,\frac{\partial \underline{u}}{\partial t}.
\]
\end{minipage}\hfill
\begin{minipage}[c]{0.36\textwidth}
\centering
\includegraphics[width=\textwidth]{infinitesimal_volume_cube.png}
\caption{Infinitesimal control volume used to relate pressure gradient to particle velocity \(\underline{u}\).}
\end{minipage}
\end{figure}

For harmonic motion, \(\underline{u}(t)=\underline{u}_0\,e^{j\omega t}\) so \(\frac{\partial}{\partial t}\to j\omega\):
\[
\frac{\partial p}{\partial x}=-j\omega \rho\,\underline{u}\qquad\Rightarrow\qquad
\underline{u}=-\frac{1}{j\omega\rho}\,\frac{\partial p}{\partial x}
= \frac{j}{k\,c\,\rho}\,\frac{\partial p}{\partial x},
\ \text{ using } k=\frac{\omega}{c}=\frac{2\pi f}{c}=\frac{2\pi}{\lambda}
\]
.

The quantity \(Z_0=c\,\rho\) is the \textbf{characteristic acoustic impedance}.

\begin{table}[H]
\centering
\caption{Wavelength \(\lambda\) and wavenumber \(k\) (assuming \(c\approx 340~\mathrm{m/s}\)).}
\begin{tabular}{c c c}
\toprule
\(f\,[\mathrm{Hz}]\) & \(\lambda\,[\mathrm{m}]\) & \(k\,[\mathrm{m^{-1}}]\)\\
\midrule
20     & 17     & 0.36 \\
200    & 3.4    & 1.34 \\
1\,000 & 0.34   & 18.47 \\
20\,000& 0.017  & 369 \\
\bottomrule
\end{tabular}
\end{table}




\subsection{Plane waves scenario}

For plane waves, the acoustic pressure can be written as the sum of forward and backward components:
\[
p(x,t)=p_{+}(t)\,e^{-jkx}+p_{-}(t)\,e^{+jkx}.
\]
In a reflection-free case (\(p_{-}=0\)):
\[
p(x,t)=p_{+}(t)\,e^{-jkx}=p_{0}\,e^{j\omega t}\,e^{-jkx}.
\]
From the previous result \(\displaystyle \underline{u}=-\frac{1}{j\omega\rho}\frac{\partial p}{\partial x}=\frac{j}{k c \rho}\frac{\partial p}{\partial x}\), we obtain
\[
\underline{u}(x,t)=\frac{j}{k c \rho}\,\frac{\partial}{\partial x}\!\left(p_{0}\,e^{j\omega t}\,e^{-jkx}\right)
=\frac{j}{k c \rho}\,(-jk)\,p_{0}\,e^{j\omega t}\,e^{-jkx}
=\frac{1}{c\rho}\,p(x,t).
\]

Hence the \textbf{specific acoustic impedance} is
\[
Z_s=\frac{p}{u}= \rho\,c \;\;\approx\; 407~\text{Rayls}\;\;(\text{air}).
\]





































 \clearpage 


 
\subsection{Summary}
\begin{tcolorbox}[colback=gray!10!white,colframe=black,title=\textbf{Summary – Electrodynamic Loudspeaker}]
\begin{itemize}
    \item A loudspeaker converts \textbf{electrical energy} into \textbf{mechanical motion} through the interaction between current and magnetic field.
    \item The Lorentz force \( \underline{F} = q\,\underline{u} \times \underline{B} \) governs the movement of the coil in the magnetic gap.
    \item The diaphragm’s oscillation generates sound waves in air.
    \item Microphones operate on the same principle, but in reverse: sound induces a voltage.
\end{itemize}
\end{tcolorbox}

\begin{tcolorbox}[colback=gray!10!white,colframe=black,title=\textbf{Summary – Force and Power in Loudspeakers}]
\begin{itemize}
    \item The total Lorentz force on the coil is $|\underline{F}| = I\,|\underline{B}|\,l\,\sin(\alpha)$.
    \item Maximum force occurs when $\underline{u} \perp \underline{B}$ ($\alpha = \frac{\pi}{2}$).
    \item Electrical power $P_E = V I$ is transformed into mechanical power $P_M = F \underline{u}$.
    \item Under ideal conditions, $P_E = P_M = P_A$ (no losses).
\end{itemize}
\end{tcolorbox}


\begin{tcolorbox}[colback=gray!10!white,colframe=black,title=\textbf{Summary – Loudspeaker Modeling Approach}]
\begin{itemize}
    \item The loudspeaker model follows the chain: Electrical $\rightarrow$ Mechanical $\rightarrow$ Acoustical.
    \item Electrical domain: $(V, I)$ — input driving voltage and current.
    \item Mechanical domain: $(\underline{F}, \underline{u})$ — force and velocity of the diaphragm.
    \item Acoustical domain: $(P, \underline{U})$ — acoustic pressure and volume velocity.
    \item The diaphragm suspension is modeled mechanically (mass, spring, and damper elements).
    \item Woofers use large cones for low frequencies; tweeters use small, light diaphragms for high frequencies.
\end{itemize}
\end{tcolorbox}


\begin{tcolorbox}[colback=gray!10!white,colframe=black,title=\textbf{Summary – Electrical and Mechanical Domains}]
\begin{itemize}
    \item \textbf{Electrical domain:}
    \begin{itemize}
        \item $V_R = R\,i$, \quad $V_C = \frac{1}{C}\int i\,dt$, \quad $V_L = L\,\frac{di}{dt}$.
        \item The three main components are resistor, capacitor, and inductor.
    \end{itemize}
    \item \textbf{Mechanical domain:}
    \begin{itemize}
        \item \textit{Mass:} $\underline{F} = M_M \frac{d\underline{u}}{dt}$ (mechanical inductance).
        \item \textit{Spring:} $\underline{u} = \frac{1}{K} \frac{d\underline{F}}{dt}$ (mechanical capacitance).
        \item \textit{Damper:} $\underline{u} = \frac{1}{R_M} \underline{F}$ (mechanical resistance).
    \end{itemize}
\end{itemize}
\end{tcolorbox}


\begin{tcolorbox}[colback=gray!10!white,colframe=black,title=\textbf{Summary – Analogy Theory}]
\begin{itemize}
    \item Two analogies are used to model electro-mechanical systems:
    \begin{itemize}
        \item \textbf{Impedance analogy (Type 1):} $\underline{F} \leftrightarrow v$, $\underline{u} \leftrightarrow i$.
        \item \textbf{Admittance analogy (Type 2):} $\underline{F} \leftrightarrow i$, $\underline{u} \leftrightarrow v$.
    \end{itemize}
    \item Type 1 is well suited for systems where force is analogous to voltage (mass modeling).
    \item Type 2 is better for systems where velocity corresponds to voltage (spring modeling).
    \item Depending on the analogy, mechanical elements correspond to different electrical components:
    \[
    \text{Mass} \leftrightarrow \text{Inductor or Capacitor}, \quad
    \text{Spring} \leftrightarrow \text{Capacitor or Inductor}, \quad
    \text{Damper} \leftrightarrow \text{Resistor or Conductor}.
    \]
\end{itemize}
\end{tcolorbox}


\begin{tcolorbox}[colback=gray!10!white,colframe=black,title=\textbf{Summary – Mechanical Element Analogs}]
\begin{itemize}
    \item \textbf{Spring:}
    \begin{itemize}
        \item Impedance analogy: capacitor $C_E$ such that ${F}(t) = \frac{1}{C_E}\int {u}(t)\,dt$.
        \item Admittance analogy: inductor $L_E$ such that ${u}(t) = L_E \frac{d{F}}{dt}$.
    \end{itemize}
    \item \textbf{Damper:}
    \begin{itemize}
        \item Impedance analogy: resistor $R_E$ such that ${u} = \frac{1}{R_E}{F}$.
        \item Admittance analogy: conductor $G_E = 1/R_E$ such that $v = \frac{1}{G_E} i$.
    \end{itemize}
    \item The choice between impedance and admittance analogy changes which electrical element represents a given mechanical component but preserves energy equivalence.
\end{itemize}
\end{tcolorbox}


\begin{tcolorbox}[colback=gray!10!white,colframe=black,title=\textbf{Summary – Mechanical Impedance}]
\begin{itemize}
    \item The power equivalence $P_M = F \, u$ corresponds to $P_E = v \, i$.
    \item Mass behaves as an inductor in the impedance analogy: $F = M_M \frac{du}{dt}$.
    \item The mechanical impedance is defined as $Z_M = \frac{F}{u}$.
    \item Impedance analogy: “across” $\leftrightarrow$ “through”.  
          Admittance analogy: “through” $\leftrightarrow$ “across”.
\end{itemize}
\end{tcolorbox}

\begin{tcolorbox}[colback=gray!10!white,colframe=black,title=\textbf{Summary – Mechanical Impedance Combinations}]
\begin{itemize}
    \item In the \textbf{impedance analogy}, mechanical impedances combine in parallel:
    \[
    Z_M = \frac{Z_{M1} Z_{M2}}{Z_{M1} + Z_{M2}}
    \]
    \item In the \textbf{admittance analogy}, mechanical impedances add directly:
    \[
    Z_M = Z_{M1} + Z_{M2}
    \]
    \item This difference arises because force and velocity swap their roles between “across” and “through” variables in each analogy.
\end{itemize}
\end{tcolorbox}


\begin{tcolorbox}[colback=gray!10!white,colframe=black,title=\textbf{Summary – Acoustic Mass}]
\begin{itemize}
    \item The acoustic domain uses \textbf{pressure} \( P \) and \textbf{volume velocity} \( U \).
    \item The acoustic mass \( M_A = \frac{M_M}{S^2} \) models the inertia of air in motion.
    \item Impedance analogy: \( P \leftrightarrow v \), \( U \leftrightarrow i \Rightarrow \) acoustic mass $\leftrightarrow$ inductor.
    \item Admittance analogy: \( P \leftrightarrow i \), \( U \leftrightarrow v \Rightarrow \) acoustic mass $\leftrightarrow$ capacitor.
\end{itemize}
\end{tcolorbox}


\begin{tcolorbox}[colback=gray!10!white,colframe=black,title=\textbf{Summary – Acoustic Resistance}]
\begin{itemize}
    \item The acoustic resistance \( R_A = \frac{R_M}{S^2} \) models energy losses due to viscosity.
    \item Acoustic resistance and mass often occur in series when modeling perforated or duct-like structures.
    \item The parameters depend on hole geometry (\( a, b, t \)) and on the properties of air (\( \rho_0, \mu \)).
    \item The total impedance of each hole is \( Z_A = R_A + Z_{M_A} \).
\end{itemize}
\end{tcolorbox}


\begin{tcolorbox}[colback=gray!10!white,colframe=black,title=\textbf{Summary – Acoustic Compliance}]
\begin{itemize}
    \item Acoustic compliance \( C_A = C_M S^2 \) quantifies the compressibility of air in a cavity.
    \item It plays a role analogous to a spring in mechanics and a capacitor in electricity.
    \item Impedance analogy: \( P \leftrightarrow v, \, U \leftrightarrow i \Rightarrow \) capacitor model.
    \item Admittance analogy: \( P \leftrightarrow i, \, U \leftrightarrow v \Rightarrow \) inductor model.
\end{itemize}
\end{tcolorbox}


\begin{tcolorbox}[colback=gray!10!white,colframe=black,title=\textbf{Summary – Specific Acoustic Impedance}]
\begin{itemize}
    \item Acoustic impedance: \( Z_A = \frac{P}{U} \) \([\frac{N \cdot s}{m^5}]\)
    \item Specific acoustic impedance: \( Z_s = \frac{P}{u} \) \([\frac{N \cdot s}{m^3}] = \text{Rayl}\)
    \item Relationship: \( Z_s = Z_A \, S \)
    \item \( Z_s \) considers molecular (particle) velocity rather than volume velocity.
\end{itemize}
\end{tcolorbox}


\begin{tcolorbox}[colback=gray!10!white,colframe=black,title=\textbf{Summary – Signal Transformer}]
\begin{itemize}
    \item A transformer converts a pair \((v, i)\) into another pair \((v', i')\) while keeping constant power \( P_E = v \, i \).
    \item The voltage–turn relationship is \( \frac{v_p}{v_s} = \frac{n_p}{n_s} = k \).
    \item The current–turn relationship is \( \frac{i_s}{i_p} = k \).
    \item Step-up transformer: \( n_s > n_p \Rightarrow v_s > v_p, \, i_s < i_p \).
    \item Step-down transformer: \( n_s < n_p \Rightarrow v_s < v_p, \, i_s > i_p \).
    \item Transformers are passive components; they transfer but do not amplify energy.
\end{itemize}
\end{tcolorbox}



\begin{tcolorbox}[colback=gray!10!white,colframe=black,title=\textbf{Remarks}]
\begin{itemize}
    \item A transformer modifies the load impedance as \( Z_p = k^2 Z_{LS} \).
    \item It allows impedance matching between circuits of different voltage and current ranges.
    \item In valve amplifiers, transformers are used to adapt the high internal impedance of the tubes to the low impedance of loudspeakers.
    \item Typical vacuum tubes provide small current (around 60 mA), hence the need for impedance transformation.
\end{itemize}
\end{tcolorbox}


\begin{tcolorbox}[colback=gray!10!white,colframe=black,title=\textbf{Summary – Gyrator Behavior}]
\begin{itemize}
    \item The gyrator maintains constant power transfer between ports:
    \[
    v_s i_s = v_p i_p
    \]
    \item The impedance transformation follows:
    \[
    Z_p = \frac{1}{g^2 Z_{LS}}
    \]
    \item If the load is:
    \begin{itemize}
        \item a resistor: \( Z_p = \frac{1}{g^2 R} \),
        \item an inductor: \( C_{eq} = g^2 L \),
        \item a capacitor: \( L_{eq} = \frac{1}{g^2 C} \).
    \end{itemize}
    \item The gyrator converts inductive behavior into capacitive behavior, and vice versa.
\end{itemize}
\end{tcolorbox}



\begin{tcolorbox}[colback=gray!10!white,colframe=black,title=\textbf{Summary – Electrical to Mechanical Transduction}]
\begin{itemize}
    \item \textbf{Admittance analogy:} represented by a transformer.
    \item \textbf{Impedance analogy:} represented by a gyrator.
    \item Relations: \( F_c = B l i \), \quad \( u_c = \dfrac{1}{B l} v \).
    \item Transconductance of the gyrator: \( g = \dfrac{1}{B l} \).
\end{itemize}
\end{tcolorbox}

\begin{tcolorbox}[colback=gray!10!white,colframe=black,title=\textbf{Summary – Mechanical to Acoustical Transduction}]
\begin{itemize}
    \item Pressure and volume velocity are linked to force and velocity through the diaphragm surface \( S_D \):
    \[
    P = \frac{F}{S_D}, \quad U = u \, S_D.
    \]
    \item Admittance analogy: transformer model (\( P \, U = F \, u \)).
    \item Impedance analogy: gyrator model.
\end{itemize}
\end{tcolorbox}

\begin{tcolorbox}[colback=gray!10!white,colframe=black,title=\textbf{Summary – Mechanical Section of the Loudspeaker}]
\begin{itemize}
    \item The diaphragm–suspension system is modeled with three main components:
    \begin{itemize}
        \item Mass \( M_M \) (diaphragm inertia)
        \item Compliance \( C_M \) (elastic restoring force)
        \item Resistance \( R_M \) (viscous damping)
    \end{itemize}
    \item The total force on the diaphragm is:
    \[
    F = \frac{1}{C_M} \int_0^t u(\tau) \, d\tau + R_M \, u + M_M \, \frac{du}{dt}
    \]
    \item This equation represents the dynamic behavior of the moving assembly of the loudspeaker.
\end{itemize}
\end{tcolorbox}


\begin{tcolorbox}[colback=gray!10!white,colframe=black,title=\textbf{Summary – Mechanical Analogies}]
\begin{itemize}
    \item \textbf{Admittance analogy:} RLC parallel circuit  
    \( \Rightarrow \) common velocity (analogous to voltage)
    \item \textbf{Impedance analogy:} RLC series circuit  
    \( \Rightarrow \) common velocity (analogous to current)
    \item The mass, compliance, and damping of the loudspeaker suspension are modeled with equivalent electrical components.
\end{itemize}
\end{tcolorbox}

\begin{tcolorbox}[colback=gray!10!white,colframe=black,title=\textbf{Summary – Loudspeaker Equivalent Model}]
\begin{itemize}
    \item The loudspeaker combines:
    \begin{itemize}
        \item An \textbf{RL circuit} (voice coil),
        \item An \textbf{RLC circuit in parallel} (mechanical suspension, admittance analogy),
        \item A \textbf{signal transformer} for both transduction stages.
    \end{itemize}
    \item Radiation impedance \( Z_{AR} = \dfrac{\bar{P}}{\bar{U}} \) relates pressure and volume velocity.
    \item Lumped element models approximate distributed systems using RLC elements.
    \item Pressure waves at the front and back of the diaphragm are in opposition:
    \[
    Z_{AR} = Z_{AB} + Z_{AF} \quad \Rightarrow \quad Y_{AR} = \frac{Y_{AF}}{2} \text{ if } Z_{AB} = Z_{AF}.
    \]
\end{itemize}
\end{tcolorbox}


\begin{tcolorbox}[colback=gray!5,colframe=gray!50!black,title=Key Points]
\begin{itemize}
    \item The transition from acoustical to mechanical domain introduces the diaphragm area \( S_D \).
    \item The transition from mechanical to electrical domain involves the electromagnetic coupling factor \( B \, l \).
    \item The moving mass \( M_{MD} \) corresponds to a capacitance \( C_{ED} = \dfrac{M_{MD}}{(B l)^2} \).
    \item A lighter diaphragm results in a higher resonance frequency and lower equivalent capacitance.
\end{itemize}
\end{tcolorbox}


\begin{tcolorbox}[colback=gray!5,colframe=gray!50!black,title=Key Points]
\begin{itemize}
    \item The suspension’s mechanical compliance \( C_{MS} \) is electrically represented as an inductance \( L_{ES} = (B l)^2 C_{MS} \).
    \item The suspension’s mechanical losses are represented as an electrical resistance \( R_{ES} = \dfrac{(B l)^2}{R_{MS}} \).
    \item High \( R_{ES} \) implies stronger damping (lower \( R_{MS} \)), reducing unwanted oscillations.
\end{itemize}
\end{tcolorbox}


\begin{tcolorbox}[colback=gray!5,colframe=gray!50!black,title=Key Points]
\begin{itemize}
    \item The complete loudspeaker model combines suspension, diaphragm, and acoustic load effects into an equivalent RLC network.
    \item The suspension compliance corresponds to the capacitance, the moving mass to the inductance, and the mechanical losses to the resistance.
    \item The resonance frequency \( f_s = \dfrac{1}{2\pi\sqrt{L_{ES} C_{ED}}} \) defines the fundamental mechanical resonance of the loudspeaker.
\end{itemize}
\end{tcolorbox}


\begin{tcolorbox}[colback=gray!5,colframe=gray!50!black,title=Key Points]
\begin{itemize}
    \item Resonance condition: \( \omega_s = \dfrac{1}{\sqrt{L C}} \).
    \item Series LC: minimum impedance \( Z_S(\omega_s) = 0 \) → short circuit.
    \item Parallel LC: maximum impedance \( Z_P(\omega_s) \rightarrow \infty \) → open circuit.
    \item Resonant LC behavior explains the frequency response of electroacoustic transducers.
\end{itemize}
\end{tcolorbox}

\begin{tcolorbox}[colback=gray!5,colframe=gray!50!black,title=Key Points]
\begin{itemize}
    \item The resonance frequency is \( \omega_s = \dfrac{1}{\sqrt{L C}} \).
    \item The quality factor is \( Q = \dfrac{R}{L \omega_s} = \dfrac{f_s}{\text{BW}} \).
    \item At resonance, \( Z_P = R \), and the circuit behaves purely resistive.
    \item This model accurately represents the loudspeaker’s electro-mechanical resonance, governed by the mass, compliance, and damping.
\end{itemize}
\end{tcolorbox}

\begin{tcolorbox}[colback=gray!5,colframe=gray!50!black,title=Key Points]
\begin{itemize}
    \item \( Z_{EM} \) models the electro-mechanical coupling and captures the resonance of the moving system.
    \item The total impedance \( Z_E \) combines electrical, inductive, and mechanical contributions.
    \item The resonance frequency \( f_s \) and quality factor \( Q_{MS} \) characterize the loudspeaker’s dynamic behavior.
    \item The efficiency \( \varepsilon = \dfrac{P_A}{P_E} \) quantifies how effectively electrical energy is transformed into acoustic power.
\end{itemize}
\end{tcolorbox}


\begin{tcolorbox}[colback=gray!5,colframe=gray!50!black,title=Key Points]
\begin{itemize}
    \item Electrical power: \( P_E = v(t) \cdot i(t) \)
    \item Mechanical power: \( P_M = F(t) \cdot u(t) \)
    \item Acoustic power: \( P_A = p(t) \cdot U(t) \)
    \item Average power: \( \bar{P} = \dfrac{v_{RMS}^2}{R} \)
    \item Crest factor: \( CF = \dfrac{P_{peak}}{\bar{P}} \), \( CF_{dB} = 10 \log_{10}(CF) \)
    \item Sinusoidal case: \( CF = 2 \; (+3\,\mathrm{dB}) \)
    \item Generic audio: \( CF \geq 4 \; (\geq +6\,\mathrm{dB}) \)
\end{itemize}
\end{tcolorbox}

\begin{tcolorbox}[colback=gray!5,colframe=gray!50!black,title=Key Points]
\begin{itemize}
    \item Instantaneous power: \( W_E(t) = v(t) i(t) = V I [\tfrac{1}{2}\cos(\beta) - \tfrac{1}{2}\cos(2 \omega t + \beta)] \)
    \item Average (active) power: \( \bar{W_E} = V_{RMS} I_{RMS} \cos(\beta) \)
    \item For \( \beta = 0 \): power is maximized (purely resistive load).
    \item For \( \beta = \frac{\pi}{2} \): average power is zero (purely reactive load).
    \item The phase \( \beta \) controls how much of the supplied energy is effectively dissipated or stored.
\end{itemize}
\end{tcolorbox}


\begin{tcolorbox}[colback=gray!5,colframe=gray!50!black,title=Key Points]
\begin{itemize}
    \item The instantaneous power is given by \( W(t) = W_p(t) + W_q(t) \).
    \item \( W_p \): in phase with voltage (active power) — produces real energy transfer.
    \item \( W_q \): in quadrature (reactive power) — represents oscillating energy with no average transfer.
    \item The average electrical power is \( \bar{W} = V_{RMS} I_{RMS} \cos(\beta) \).
\end{itemize}
\end{tcolorbox}


\begin{tcolorbox}[colback=gray!5,colframe=gray!50!black,title=Key Points]
\begin{itemize}
    \item The current can be decomposed as \( i(t) = i_p(t) + j i_q(t) \), where:
    \[
    i_p(t) = \frac{\mathrm{Re}[Z_E]}{|Z_E|^2} v(t), \quad i_q(t) = \frac{\mathrm{Im}[Z_E]}{|Z_E|^2} v(t)
    \]
    \item The instantaneous power has both real and imaginary components:
    \[
    W(t) = W_p(t) + j W_q(t)
    \]
    \item The average power only depends on the real part of the impedance:
    \[
    \bar{W} = \mathrm{Re}[Z_E] I_{RMS}^2
    \]
    \item In a purely conservative (lossless) system, \( \mathrm{Re}[Z_E] = 0 \), and thus \( \bar{W} = 0 \).
\end{itemize}
\end{tcolorbox}

\begin{tcolorbox}[colback=gray!5,colframe=gray!50!black,title=Key results (mechanical domain)]
\[
\boxed{\,u(t)=\dfrac{F(t)}{Z_M},\quad
W_M(t)=F(t)u(t),\quad
\bar W_M=R_M\,u_{\mathrm{RMS}}^{\,2}\,}
\]
Only \(R_M\) contributes to mean power; \(X_M\) stores and returns energy with zero average transfer.
\end{tcolorbox}

\begin{tcolorbox}[colback=gray!5,colframe=gray!50!black,title=Key results (acoustical domain)]
\[
\boxed{\,U=\dfrac{P}{Z_A},\quad W_A=P\,U,\quad
\bar W_A=R_A\,U_{\mathrm{RMS}}^{\,2}
=\dfrac{R_A}{|Z_A|^2}P_{\mathrm{RMS}}^{\,2}\,}
\]
Only \(R_A\) contributes to mean radiated power; \(X_A\) stores and releases energy with zero average transfer.
\end{tcolorbox}


\begin{tcolorbox}[colback=gray!5,colframe=gray!50!black,title=Key points]
\begin{itemize}
\item Momentum balance on \(dV\): \(-\partial p/\partial x=\rho\,\partial\underline{u}/\partial t\).
\item Harmonic regime: \(\underline{u}=-\dfrac{1}{j\omega\rho}\,\dfrac{\partial p}{\partial x}=\dfrac{j}{k c \rho}\,\dfrac{\partial p}{\partial x}\).
\item Dispersion relations: \(k=\dfrac{\omega}{c}=\dfrac{2\pi f}{c}=\dfrac{2\pi}{\lambda}\).
\item Characteristic impedance: \(Z_0=c\rho\).
\end{itemize}
\end{tcolorbox}